
\documentclass{article}
\usepackage{colortbl}
\usepackage{makecell}
\usepackage{multirow}
\usepackage{supertabular}

\begin{document}

\newcounter{utterance}

\centering \large Interaction Transcript for game `clean\_up', experiment `2\_hard\_3obj\_en', episode 1 with kimi{-}k2.5{-}without{-}reasoning{-}t1.0.
\vspace{24pt}

{ \footnotesize  \setcounter{utterance}{1}
\setlength{\tabcolsep}{0pt}
\begin{supertabular}{c@{$\;$}|p{.15\linewidth}@{}p{.15\linewidth}p{.15\linewidth}p{.15\linewidth}p{.15\linewidth}p{.15\linewidth}}
    \# & $\;$A & \multicolumn{4}{c}{Game Master} & $\;\:$B\\
    \hline

    \theutterance \stepcounter{utterance}  
    & & \multicolumn{4}{p{0.6\linewidth}}{
        \cellcolor[rgb]{0.9,0.9,0.9}{
            \makecell[{{p{\linewidth}}}]{
                \texttt{\tiny{[P1$\langle$GM]}}
                \texttt{I am your game master, and you are playing a collaborative game with the following grid as game board:} \\
\\ 
\texttt{```} \\
\texttt{ 1234567} \\
\texttt{╔═══╤═╤═╗} \\
\texttt{║◌◌◌│◌│◌║ 1} \\
\texttt{╟─┬─┼─┴─╢ 2} \\
\texttt{║◌│◌│C◌◌║ 3} \\
\texttt{║◌├─┤◌┌─╢ 4} \\
\texttt{╟─┘◌│◌│L║ 5} \\
\texttt{║◌◌◌└┬┴─╢ 6} \\
\texttt{║P◌◌◌│◌◌║ 7} \\
\texttt{╚════╧══╝} \\
\\ 
\texttt{```} \\
\\ 
\texttt{* The upper edge displays x{-}coordinates increasing to the right, and the right edge y{-}coordinates increasing downward.} \\
\texttt{* The following objects are randomly placed on your grid: 'C', 'L', 'P'.} \\
\\ 
\\ 
\texttt{The other player sees a variation of the game board, where the objects are placed at different random locations. You cannot see the other player's board, and they cannot see yours.} \\
\\ 
\\ 
\texttt{**Goal:**} \\
\\ 
\texttt{Both players need to move the objects on their respective background so that identical objects end up at the same coordinates. You have to communicate with the other player to agree upon a common goal configuration.} \\
\\ 
\\ 
\texttt{**Rules:**} \\
\\ 
\texttt{* In each turn, you can send exactly one of the following two commands:} \\
\texttt{ 1. `SAY: <MESSAGE>`: to send a message (everything up until the next line break) to the other player. I will forward it to your partner.} \\
\texttt{2. `MOVE: <OBJECT>, (<X>, <Y>)`: to move an object to a new position, where `<X>` is the column and `<Y>` is the row. I will inform you if your move was successful or not.} \\
\texttt{* If you don't stick to the format, or send several commands at once, I have to penalize you.} \\
\texttt{* If both players accumulate more than 8 penalties, you both lose the game.} \\
\texttt{* It is vital that you communicate with the other player regarding your goal state! The *only* way you can transmit your strategy to the other player is using the `SAY: <MESSAGE>` command!} \\
\\ 
\\ 
\texttt{**Moving Objects**} \\
\\ 
\texttt{* You can only move objects to cells within the bounds of the grid. The target cell must be empty, i.e., it must only contain the symbol '◌'.} \\
\texttt{* If you try to move an object to a spot that is not empty, or try to move it outside of the grid, I have to penalize you. You get another try.} \\
\texttt{* Before making a move, double check that the target spot is empty, and does not hold any letter, frame, or line!} \\
\\ 
\\ 
\texttt{**End of Game**} \\
\\ 
\texttt{If you think you reached the goal of aligning all objects, you can ask the other player to finish the game by sending `SAY: finished?`. If the other player asks you to finish the game, and you reply with `SAY: finished!`, the game will end.} \\
\\ 
\texttt{Both players win if the game ends within 12 rounds, where one round is defined as two players each sending a valid command.} \\
\\ 
\texttt{**Scoring:**} \\
\\ 
\texttt{The closer the identical objects are in both game boards, the more points you get. Penalties reduce your points. Can you beat the record?} \\
\\ 
\\ 
\texttt{Please send a message to the other player to start the game!} \\
            }
        }
    }
    & \\ \\

    \theutterance \stepcounter{utterance}  
    & \multicolumn{4}{p{0.6\linewidth}}{
        \cellcolor[rgb]{0.9,0.9,0.9}{
            \makecell[{{p{\linewidth}}}]{
                \texttt{\tiny{[P1$\rangle$GM]}}
                \texttt{SAY: Hi! Let's solve this together. I can see C at (3,3), L at (7,5), and P at (1,7) on my board. Where are your objects located? Once we know both configurations, we can agree on a common target position for each object.} \\
            }
        }
    }
    & & \\ \\

    \theutterance \stepcounter{utterance}  
    & & & \multicolumn{2}{p{0.3\linewidth}}{
        \cellcolor[rgb]{0.9,0.9,0.9}{
            \makecell[{{p{\linewidth}}}]{
                \texttt{\tiny{[GM$|$GM]}}
                \texttt{SAY: Hi! Let's solve this together. I can see C at (3,3), L at (7,5), and P at (1,7) on my board. Where are your objects located? Once we know both configurations, we can agree on a common target position for each object.} \\
            }
        }
    }
    & & \\ \\

    \theutterance \stepcounter{utterance}  
    & & \multicolumn{4}{p{0.6\linewidth}}{
        \cellcolor[rgb]{0.9,0.9,0.9}{
            \makecell[{{p{\linewidth}}}]{
                \texttt{\tiny{[GM$\rangle$P2]}}
                \texttt{I am your game master, and you are playing a collaborative game with the following grid as game board:} \\
\\ 
\texttt{```} \\
\texttt{ 1234567} \\
\texttt{╔═══╤═╤═╗} \\
\texttt{║◌◌◌│◌│◌║ 1} \\
\texttt{╟─┬─┼─┴─╢ 2} \\
\texttt{║◌│L│◌P◌║ 3} \\
\texttt{║◌├─┤◌┌─╢ 4} \\
\texttt{╟─┘◌│◌│◌║ 5} \\
\texttt{║◌◌◌└┬┴─╢ 6} \\
\texttt{║◌C◌◌│◌◌║ 7} \\
\texttt{╚════╧══╝} \\
\\ 
\texttt{```} \\
\\ 
\texttt{* The upper edge displays x{-}coordinates increasing to the right, and the right edge y{-}coordinates increasing downward.} \\
\texttt{* The following objects are randomly placed on your grid: 'C', 'L', 'P'.} \\
\\ 
\\ 
\texttt{The other player sees a variation of the game board, where the objects are placed at different random locations. You cannot see the other player's board, and they cannot see yours.} \\
\\ 
\\ 
\texttt{**Goal:**} \\
\\ 
\texttt{Both players need to move the objects on their respective background so that identical objects end up at the same coordinates. You have to communicate with the other player to agree upon a common goal configuration.} \\
\\ 
\\ 
\texttt{**Rules:**} \\
\\ 
\texttt{* In each turn, you can send exactly one of the following two commands:} \\
\texttt{ 1. `SAY: <MESSAGE>`: to send a message (everything up until the next line break) to the other player. I will forward it to your partner.} \\
\texttt{2. `MOVE: <OBJECT>, (<X>, <Y>)`: to move an object to a new position, where `<X>` is the column and `<Y>` is the row. I will inform you if your move was successful or not.} \\
\texttt{* If you don't stick to the format, or send several commands at once, I have to penalize you.} \\
\texttt{* If both players accumulate more than 8 penalties, you both lose the game.} \\
\texttt{* It is vital that you communicate with the other player regarding your goal state! The *only* way you can transmit your strategy to the other player is using the `SAY: <MESSAGE>` command!} \\
\\ 
\\ 
\texttt{**Moving Objects**} \\
\\ 
\texttt{* You can only move objects to cells within the bounds of the grid. The target cell must be empty, i.e., it must only contain the symbol '◌'.} \\
\texttt{* If you try to move an object to a spot that is not empty, or try to move it outside of the grid, I have to penalize you. You get another try.} \\
\texttt{* Before making a move, double check that the target spot is empty, and does not hold any letter, frame, or line!} \\
\\ 
\\ 
\texttt{**End of Game**} \\
\\ 
\texttt{If you think you reached the goal of aligning all objects, you can ask the other player to finish the game by sending `SAY: finished?`. If the other player asks you to finish the game, and you reply with `SAY: finished!`, the game will end.} \\
\\ 
\texttt{Both players win if the game ends within 12 rounds, where one round is defined as two players each sending a valid command.} \\
\\ 
\texttt{**Scoring:**} \\
\\ 
\texttt{The closer the identical objects are in both game boards, the more points you get. Penalties reduce your points. Can you beat the record?} \\
\\ 
\texttt{The other player started the game by sending this message:} \\
\\ 
\texttt{"Hi! Let's solve this together. I can see C at (3,3), L at (7,5), and P at (1,7) on my board. Where are your objects located? Once we know both configurations, we can agree on a common target position for each object."} \\
\\ 
\texttt{What is your first command?} \\
            }
        }
    }
    & \\ \\

    \theutterance \stepcounter{utterance}  
    & & & \multicolumn{4}{p{0.6\linewidth}}{
        \cellcolor[rgb]{0.9,0.9,0.9}{
            \makecell[{{p{\linewidth}}}]{
                \texttt{\tiny{[GM$\langle$P2]}}
                \texttt{SAY: Hi! On my board, I can see C at (2,7), L at (3,3), and P at (5,3). So we have:} \\
\texttt{{-} C: I'm at (2,7), you're at (3,3)} \\
\texttt{{-} L: I'm at (3,3), you're at (7,5)} \\
\texttt{{-} P: I'm at (5,3), you're at (1,7)} \\
\\ 
\texttt{Let me suggest we move all objects to the center area. How about we aim for:} \\
\texttt{{-} C at (4,4)} \\
\texttt{{-} L at (4,5)} \\
\texttt{{-} P at (5,4)} \\
\\ 
\texttt{Or if you have a different preference, let me know!} \\
            }
        }
    }
     \\ \\

    \theutterance \stepcounter{utterance}  
    & & & \multicolumn{2}{p{0.3\linewidth}}{
        \cellcolor[rgb]{0.9,0.9,0.9}{
            \makecell[{{p{\linewidth}}}]{
                \texttt{\tiny{[GM$|$GM]}}
                \texttt{SAY: Hi! On my board, I can see C at (2,7), L at (3,3), and P at (5,3). So we have:} \\
\texttt{{-} C: I'm at (2,7), you're at (3,3)} \\
\texttt{{-} L: I'm at (3,3), you're at (7,5)} \\
\texttt{{-} P: I'm at (5,3), you're at (1,7)} \\
\\ 
\texttt{Let me suggest we move all objects to the center area. How about we aim for:} \\
\texttt{{-} C at (4,4)} \\
\texttt{{-} L at (4,5)} \\
\texttt{{-} P at (5,4)} \\
\\ 
\texttt{Or if you have a different preference, let me know!} \\
            }
        }
    }
    & & \\ \\

    \theutterance \stepcounter{utterance}  
    & & & \multicolumn{2}{p{0.3\linewidth}}{
        \cellcolor[rgb]{0.9,0.9,0.9}{
            \makecell[{{p{\linewidth}}}]{
                \texttt{\tiny{[GM$|$GM]}}
                \texttt{Invalid format: tail is not empty: '} \\
\texttt{{-} C: I'm at (2,7), you're at (3,3)} \\
\texttt{{-} L: I'm at (3,3), you're at (7,5)} \\
\texttt{{-} P: I'm at (5,3), you're at (1,7)} \\
\\ 
\texttt{Let me suggest we move all objects to the center area. How about we aim for:} \\
\texttt{{-} C at (4,4)} \\
\texttt{{-} L at (4,5)} \\
\texttt{{-} P at (5,4)} \\
\\ 
\texttt{Or if you have a different preference, let me know!} \\
            }
        }
    }
    & & \\ \\

    \theutterance \stepcounter{utterance}  
    & & \multicolumn{4}{p{0.6\linewidth}}{
        \cellcolor[rgb]{0.9,0.9,0.9}{
            \makecell[{{p{\linewidth}}}]{
                \texttt{\tiny{[GM$\rangle$P2]}}
                \texttt{Penalty: Your message must not contain anything after the command!} \\
\\ 
\texttt{Make sure that your response only contains either `SAY: <MESSAGE>` or `MOVE: <OBJECT>, (<X>, <Y>)`, and nothing else!} \\
\\ 
\texttt{You have collectively accumulated 1 of 8 penalties.} \\
\texttt{Please try again!} \\
            }
        }
    }
    & \\ \\

    \theutterance \stepcounter{utterance}  
    & & & \multicolumn{4}{p{0.6\linewidth}}{
        \cellcolor[rgb]{0.9,0.9,0.9}{
            \makecell[{{p{\linewidth}}}]{
                \texttt{\tiny{[GM$\langle$P2]}}
                \texttt{SAY: Hi! On my board, C is at (2,7), L is at (3,3), and P is at (5,3). I suggest we move all objects to the center: C to (4,4), L to (4,5), P to (5,4). What do you think?} \\
            }
        }
    }
     \\ \\

    \theutterance \stepcounter{utterance}  
    & & & \multicolumn{2}{p{0.3\linewidth}}{
        \cellcolor[rgb]{0.9,0.9,0.9}{
            \makecell[{{p{\linewidth}}}]{
                \texttt{\tiny{[GM$|$GM]}}
                \texttt{SAY: Hi! On my board, C is at (2,7), L is at (3,3), and P is at (5,3). I suggest we move all objects to the center: C to (4,4), L to (4,5), P to (5,4). What do you think?} \\
            }
        }
    }
    & & \\ \\

    \theutterance \stepcounter{utterance}  
    & & \multicolumn{4}{p{0.6\linewidth}}{
        \cellcolor[rgb]{0.9,0.9,0.9}{
            \makecell[{{p{\linewidth}}}]{
                \texttt{\tiny{[P1$\langle$GM]}}
                \texttt{Your message has been relayed to the other player.} \\
\\ 
\\ 
\texttt{You are currently playing round 1 of maximum 12.} \\
\texttt{You have collectively accumulated 1 of 8 penalties.} \\
\texttt{The other player sent this message:} \\
\\ 
\\ 
\texttt{"Hi! On my board, C is at (2,7), L is at (3,3), and P is at (5,3). I suggest we move all objects to the center: C to (4,4), L to (4,5), P to (5,4). What do you think?"} \\
\\ 
\\ 
\texttt{What is your next command?} \\
            }
        }
    }
    & \\ \\

    \theutterance \stepcounter{utterance}  
    & \multicolumn{4}{p{0.6\linewidth}}{
        \cellcolor[rgb]{0.9,0.9,0.9}{
            \makecell[{{p{\linewidth}}}]{
                \texttt{\tiny{[P1$\rangle$GM]}}
                \texttt{SAY: I like the idea of meeting in the center! However, I need to check if those positions are empty on my board. Let me verify: (4,4) has '◌', (4,5) has '◌', (5,4) has '◌' {-} all clear! Your plan works for me. Let's do it: C→(4,4), L→(4,5), P→(5,4). I'll start moving. Should we agree on an order {-} maybe move C first, then L, then P?} \\
            }
        }
    }
    & & \\ \\

    \theutterance \stepcounter{utterance}  
    & & & \multicolumn{2}{p{0.3\linewidth}}{
        \cellcolor[rgb]{0.9,0.9,0.9}{
            \makecell[{{p{\linewidth}}}]{
                \texttt{\tiny{[GM$|$GM]}}
                \texttt{SAY: I like the idea of meeting in the center! However, I need to check if those positions are empty on my board. Let me verify: (4,4) has '◌', (4,5) has '◌', (5,4) has '◌' {-} all clear! Your plan works for me. Let's do it: C→(4,4), L→(4,5), P→(5,4). I'll start moving. Should we agree on an order {-} maybe move C first, then L, then P?} \\
            }
        }
    }
    & & \\ \\

    \theutterance \stepcounter{utterance}  
    & & \multicolumn{4}{p{0.6\linewidth}}{
        \cellcolor[rgb]{0.9,0.9,0.9}{
            \makecell[{{p{\linewidth}}}]{
                \texttt{\tiny{[GM$\rangle$P2]}}
                \texttt{Your message has been relayed to the other player.} \\
\\ 
\\ 
\texttt{You are currently playing round 2 of maximum 12.} \\
\texttt{You have collectively accumulated 1 of 8 penalties.} \\
\texttt{The other player sent this message:} \\
\\ 
\\ 
\texttt{"I like the idea of meeting in the center! However, I need to check if those positions are empty on my board. Let me verify: (4,4) has '◌', (4,5) has '◌', (5,4) has '◌' {-} all clear! Your plan works for me. Let's do it: C→(4,4), L→(4,5), P→(5,4). I'll start moving. Should we agree on an order {-} maybe move C first, then L, then P?"} \\
\\ 
\\ 
\texttt{What is your next command?} \\
            }
        }
    }
    & \\ \\

    \theutterance \stepcounter{utterance}  
    & & & \multicolumn{4}{p{0.6\linewidth}}{
        \cellcolor[rgb]{0.9,0.9,0.9}{
            \makecell[{{p{\linewidth}}}]{
                \texttt{\tiny{[GM$\langle$P2]}}
                \texttt{SAY: Sounds good! Let's move C first, then L, then P. I'll start by moving my C from (2,7) to (4,4).} \\
            }
        }
    }
     \\ \\

    \theutterance \stepcounter{utterance}  
    & & & \multicolumn{2}{p{0.3\linewidth}}{
        \cellcolor[rgb]{0.9,0.9,0.9}{
            \makecell[{{p{\linewidth}}}]{
                \texttt{\tiny{[GM$|$GM]}}
                \texttt{SAY: Sounds good! Let's move C first, then L, then P. I'll start by moving my C from (2,7) to (4,4).} \\
            }
        }
    }
    & & \\ \\

    \theutterance \stepcounter{utterance}  
    & & \multicolumn{4}{p{0.6\linewidth}}{
        \cellcolor[rgb]{0.9,0.9,0.9}{
            \makecell[{{p{\linewidth}}}]{
                \texttt{\tiny{[P1$\langle$GM]}}
                \texttt{Your message has been relayed to the other player.} \\
\\ 
\\ 
\texttt{You are currently playing round 2 of maximum 12.} \\
\texttt{You have collectively accumulated 1 of 8 penalties.} \\
\texttt{The other player sent this message:} \\
\\ 
\\ 
\texttt{"Sounds good! Let's move C first, then L, then P. I'll start by moving my C from (2,7) to (4,4)."} \\
\\ 
\\ 
\texttt{What is your next command?} \\
            }
        }
    }
    & \\ \\

    \theutterance \stepcounter{utterance}  
    & \multicolumn{4}{p{0.6\linewidth}}{
        \cellcolor[rgb]{0.9,0.9,0.9}{
            \makecell[{{p{\linewidth}}}]{
                \texttt{\tiny{[P1$\rangle$GM]}}
                \texttt{MOVE: C, (4, 4)} \\
            }
        }
    }
    & & \\ \\

    \theutterance \stepcounter{utterance}  
    & & & \multicolumn{2}{p{0.3\linewidth}}{
        \cellcolor[rgb]{0.9,0.9,0.9}{
            \makecell[{{p{\linewidth}}}]{
                \texttt{\tiny{[GM$|$GM]}}
                \texttt{MOVE: C, (4, 4)} \\
            }
        }
    }
    & & \\ \\

    \theutterance \stepcounter{utterance}  
    & & & \multicolumn{2}{p{0.3\linewidth}}{
        \cellcolor[rgb]{0.9,0.9,0.9}{
            \makecell[{{p{\linewidth}}}]{
                \texttt{\tiny{[GM$|$GM]}}
                \texttt{Penalty: (4,4) is not empty, but contains '┤'.} \\
\\ 
\texttt{You have collectively accumulated 2 of 8 penalties.} \\
\texttt{Please try again!} \\
            }
        }
    }
    & & \\ \\

    \theutterance \stepcounter{utterance}  
    & & \multicolumn{4}{p{0.6\linewidth}}{
        \cellcolor[rgb]{0.9,0.9,0.9}{
            \makecell[{{p{\linewidth}}}]{
                \texttt{\tiny{[P1$\langle$GM]}}
                \texttt{Penalty: (4,4) is not empty, but contains '┤'.} \\
\\ 
\texttt{You have collectively accumulated 2 of 8 penalties.} \\
\texttt{Please try again!} \\
            }
        }
    }
    & \\ \\

    \theutterance \stepcounter{utterance}  
    & \multicolumn{4}{p{0.6\linewidth}}{
        \cellcolor[rgb]{0.9,0.9,0.9}{
            \makecell[{{p{\linewidth}}}]{
                \texttt{\tiny{[P1$\rangle$GM]}}
                \texttt{SAY: Oops! I got a penalty {-} (4,4) has '┤' on my board, not empty. Let me re{-}check my board more carefully. Looking again: C is at (3,3), L at (7,5), P at (1,7). The center area has walls. Can you tell me which positions are empty on your board? Then I can suggest alternative target coordinates that work for both of us.} \\
            }
        }
    }
    & & \\ \\

    \theutterance \stepcounter{utterance}  
    & & & \multicolumn{2}{p{0.3\linewidth}}{
        \cellcolor[rgb]{0.9,0.9,0.9}{
            \makecell[{{p{\linewidth}}}]{
                \texttt{\tiny{[GM$|$GM]}}
                \texttt{SAY: Oops! I got a penalty {-} (4,4) has '┤' on my board, not empty. Let me re{-}check my board more carefully. Looking again: C is at (3,3), L at (7,5), P at (1,7). The center area has walls. Can you tell me which positions are empty on your board? Then I can suggest alternative target coordinates that work for both of us.} \\
            }
        }
    }
    & & \\ \\

    \theutterance \stepcounter{utterance}  
    & & \multicolumn{4}{p{0.6\linewidth}}{
        \cellcolor[rgb]{0.9,0.9,0.9}{
            \makecell[{{p{\linewidth}}}]{
                \texttt{\tiny{[GM$\rangle$P2]}}
                \texttt{Your message has been relayed to the other player.} \\
\\ 
\\ 
\texttt{You are currently playing round 3 of maximum 12.} \\
\texttt{You have collectively accumulated 2 of 8 penalties.} \\
\texttt{The other player sent this message:} \\
\\ 
\\ 
\texttt{"Oops! I got a penalty {-} (4,4) has '┤' on my board, not empty. Let me re{-}check my board more carefully. Looking again: C is at (3,3), L at (7,5), P at (1,7). The center area has walls. Can you tell me which positions are empty on your board? Then I can suggest alternative target coordinates that work for both of us."} \\
\\ 
\\ 
\texttt{What is your next command?} \\
            }
        }
    }
    & \\ \\

    \theutterance \stepcounter{utterance}  
    & & & \multicolumn{4}{p{0.6\linewidth}}{
        \cellcolor[rgb]{0.9,0.9,0.9}{
            \makecell[{{p{\linewidth}}}]{
                \texttt{\tiny{[GM$\langle$P2]}}
                \texttt{SAY: On my board, empty cells include: (4,4), (4,5), (5,4), (6,4), (6,5), (2,2), (2,6), (6,2), (6,6). Let me know which of these work for you, or suggest others!} \\
            }
        }
    }
     \\ \\

    \theutterance \stepcounter{utterance}  
    & & & \multicolumn{2}{p{0.3\linewidth}}{
        \cellcolor[rgb]{0.9,0.9,0.9}{
            \makecell[{{p{\linewidth}}}]{
                \texttt{\tiny{[GM$|$GM]}}
                \texttt{SAY: On my board, empty cells include: (4,4), (4,5), (5,4), (6,4), (6,5), (2,2), (2,6), (6,2), (6,6). Let me know which of these work for you, or suggest others!} \\
            }
        }
    }
    & & \\ \\

    \theutterance \stepcounter{utterance}  
    & & \multicolumn{4}{p{0.6\linewidth}}{
        \cellcolor[rgb]{0.9,0.9,0.9}{
            \makecell[{{p{\linewidth}}}]{
                \texttt{\tiny{[P1$\langle$GM]}}
                \texttt{Your message has been relayed to the other player.} \\
\\ 
\\ 
\texttt{You are currently playing round 3 of maximum 12.} \\
\texttt{You have collectively accumulated 2 of 8 penalties.} \\
\texttt{The other player sent this message:} \\
\\ 
\\ 
\texttt{"On my board, empty cells include: (4,4), (4,5), (5,4), (6,4), (6,5), (2,2), (2,6), (6,2), (6,6). Let me know which of these work for you, or suggest others!"} \\
\\ 
\\ 
\texttt{What is your next command?} \\
            }
        }
    }
    & \\ \\

    \theutterance \stepcounter{utterance}  
    & \multicolumn{4}{p{0.6\linewidth}}{
        \cellcolor[rgb]{0.9,0.9,0.9}{
            \makecell[{{p{\linewidth}}}]{
                \texttt{\tiny{[P1$\rangle$GM]}}
                \texttt{Let me carefully check my board for empty cells. Looking at the grid:} \\
\\ 
\texttt{```} \\
\texttt{ 1234567} \\
\texttt{╔═══╤═╤═╗} \\
\texttt{║◌◌◌│◌│◌║ 1} \\
\texttt{╟─┬─┼─┴─╢ 2} \\
\texttt{║◌│◌│C◌◌║ 3} \\
\texttt{║◌├─┤◌┌─╢ 4} \\
\texttt{╟─┘◌│◌│L║ 5} \\
\texttt{║◌◌◌└┬┴─╢ 6} \\
\texttt{║P◌◌◌│◌◌║ 7} \\
\texttt{╚════╧══╝} \\
\texttt{```} \\
\\ 
\texttt{Let me identify empty '◌' cells by coordinates:} \\
\texttt{{-} Row 1: (1,1), (2,1), (3,1), (5,1), (7,1)} \\
\texttt{{-} Row 2: (1,2), (3,2), (5,2), (6,2), (7,2)} \\
\texttt{{-} Row 3: (1,3), (2,3), (4,3), (5,3), (6,3), (7,3) {-} but C is at (3,3)} \\
\texttt{{-} Row 4: (1,4), (4,4) has '┤', (5,4), (6,4) {-} wait let me re{-}read} \\
\\ 
\texttt{Actually looking more carefully at row 4: `║◌├─┤◌┌─╢` {-} positions are: 1=◌, 2=├, 3=─, 4=┤, 5=◌, 6=┌, 7=─} \\
\\ 
\texttt{So (4,4) is '┤', (5,4) is '◌', (6,4) is '┌'} \\
\\ 
\texttt{Row 5: `╟─┘◌│◌│L║` {-} (3,5)=◌, (5,5)=◌, and L at (7,5)} \\
\\ 
\texttt{Row 6: `║◌◌◌└┬┴─╢` {-} (1,6)=◌, (2,6)=◌, (3,6)=◌, (4,6)=└, (5,6)=┬, (6,6)=┴} \\
\\ 
\texttt{Row 7: `║P◌◌◌│◌◌║` {-} P at (1,7), (2,7)=◌, (3,7)=◌, (4,7)=◌, (6,7)=◌, (7,7)=◌} \\
\\ 
\texttt{From the other player's list, checking my board: (6,4) has '┌', (6,5) has '│', (6,6) has '┴'} \\
\\ 
\texttt{Common empty cells: (2,2), (2,6), (6,2) {-} let me verify (6,2): row 2 `╟─┬─┼─┴─╢` {-} position 6 is '┴', not ◌.} \\
\\ 
\texttt{Wait let me re{-}check. Row 2: `╟─┬─┼─┴─╢` {-} this is a border row. The characters are: ╟, ─, ┬, ─, ┼, ─, ┴, ─, ╢} \\
\\ 
\texttt{So positions: 1=╟, 2=─, 3=┬, 4=─, 5=┼, 6=─, 7=┴, 8=─, 9=╢? That doesn't match 7 columns.} \\
\\ 
\texttt{Let me re{-}interpret: the grid shows the internal cells. The borders are separate.} \\
\\ 
\texttt{Looking again at row 2: `╟─┬─┼─┴─╢` {-} this is the border between rows 1 and 2. The actual row 2 cells are above this.} \\
\\ 
\texttt{Actually, re{-}reading: row 3 shows `║◌│◌│C◌◌║` {-} so the pattern is: ║, then cells, then ║} \\
\\ 
\texttt{For row 2, the border `╟─┬─┼─┴─╢` indicates connections. The cells in row 2 would be at positions where there are vertical connections.} \\
\\ 
\texttt{Let me look at row 3 again: `║◌│◌│C◌◌║` {-} positions 1,2,3,4,5,6,7: ◌, │, ◌, │, C, ◌, ◌} \\
\\ 
\texttt{So (1,3)=◌, (2,3)=│, (3,3)=C, wait that's wrong. Let me recount: `║◌│◌│C◌◌║` has characters: ║, ◌, │, ◌, │, C, ◌, ◌, ║} \\
\\ 
\texttt{So positions: 1=◌, 2=│, 3=◌, 4=│, 5=C, 6=◌, 7=◌? That's 7 positions plus borders.} \\
\\ 
\texttt{Hmm, but C is shown in the middle. Let me try: ◌, │, ◌, │, C, ◌, ◌ {-} that's 7 chars between ║ and ║.} \\
\\ 
\texttt{So (1,3)=◌, (2,3)=│, (3,3)=◌, (4,3)=│, (5,3)=C? No wait, the C is at position 3 in the display `C◌◌`.} \\
\\ 
\texttt{Actually `│C◌◌` would be positions 2,3,4,5. So (2,3)=│, (3,3)=C, (4,3)=◌, (5,3)=◌? But then what about (1,3)?} \\
\\ 
\texttt{Let me try: `◌│◌│C◌◌` = positions 1{-}7: (1,3)=◌, (2,3)=│, (3,3)=◌, (4,3)=│, (5,3)=C, (6,3)=◌, (7,3)=◌} \\
\\ 
\texttt{But the string is `◌│◌│C◌◌` which is 7 chars. Yes! So C is at (5,3)? But I was told C is at (3,3)...} \\
\\ 
\texttt{Wait, the problem says C is at (3,3). Let me re{-}read: `║◌│◌│C◌◌║` {-} maybe the │ are cell separators, not cell contents?} \\
\\ 
\texttt{Actually looking at the pattern, I think the grid uses box{-}drawing characters where │, ┤, ├, etc. are walls/part of the maze structure, and ◌ represents empty walkable cells.} \\
\\ 
\texttt{Let me re{-}interpret: in `║◌│◌│C◌◌║`, the walkable cells are where ◌ and letters are. The │ are vertical walls.} \\
\\ 
\texttt{So positions with walkable cells: 1=◌, 3=◌, 5=C, 6=◌, 7=◌? That's only 5 positions for 7 columns.} \\
\\ 
\texttt{Hmm, this is confusing. Let me try a different approach {-} look at row 1: `║◌◌◌│◌│◌║` {-} this has ◌◌◌│◌│◌ between the ║. That's 7 characters: ◌,◌,◌,│,◌,│,◌} \\
\\ 
\texttt{So columns 1,2,3 are ◌, column 4 is │, column 5 is ◌, column 6 is │, column 7 is ◌.} \\
\\ 
\texttt{For row 3: `║◌│◌│C◌◌║` = ◌,│,◌,│,C,◌,◌ so: (1,3)=◌, (2,3)=│, (3,3)=◌, (4,3)=│, (5,3)=C, (6,3)=◌, (7,3)=◌} \\
\\ 
\texttt{But the problem states C is at (3,3), not (5,3). } \\
\\ 
\texttt{Wait {-} re{-}reading the problem: "The following objects are randomly placed on your grid: 'C', 'L', 'P'." and looking at the board, C appears in the middle of row 3. If columns are 1{-}7 left to right, and C is the 5th visible character in `◌│◌│C◌◌`, that would be column 5.} \\
\\ 
\texttt{Unless... the coordinates refer to the ◌ positions only? No, that doesn't make sense.} \\
\\ 
\texttt{Let me re{-}check: maybe I miscounted. `║◌│◌│C◌◌║` {-} between ║ and ║ we have: ◌, │, ◌, │, C, ◌, ◌. That's 7 characters for 7 columns. C is at position 5.} \\
\\ 
\texttt{But the problem says "C at (3,3)" in my initial understanding. Actually, re{-}reading my own message {-} I said "I can see C at (3,3)" {-} I assumed that based on visual center, but let me verify by counting.} \\
\\ 
\texttt{Actually wait {-} I need to re{-}read more carefully. The user said I claimed "C at (3,3)" in my first message. Let me check if that's correct by looking at the string again.} \\
\\ 
\texttt{Row 3: `║◌│◌│C◌◌║` {-} positions: 1=◌, 2=│, 3=◌, 4=│, 5=C, 6=◌, 7=◌} \\
\\ 
\texttt{Hmm, but I said (3,3). Unless the │ characters don't count as columns? That would be strange.} \\
\\ 
\texttt{Actually, re{-}looking: maybe the grid is showing that some columns are narrower? Or perhaps I'm misunderstanding the encoding.} \\
\\ 
\texttt{Let me try yet another interpretation: the characters like │, ┤, ├, ─, etc. are part of the box{-}drawing and the actual cell contents are only the letters and ◌.} \\
\\ 
\texttt{In row 3: between the vertical borders, we have cells separated by vertical lines. So `◌│◌│C◌◌` could mean: cell1=◌, wall, cell2=◌, wall, cell3=C, cell4=◌, cell5=◌? That would be 5 cells.} \\
\\ 
\texttt{But we need 7 columns. } \\
\\ 
\texttt{Let me look at row 1 again: `║◌◌◌│◌│◌║` {-} if this is 7 cells, then │ is a cell content, not a separator.} \\
\\ 
\texttt{I think I need to accept that: columns 1{-}7 contain the characters as shown, and │, ┤, etc. are obstacles/walls that block movement. The ◌ and letters are walkable.} \\
\\ 
\texttt{So for row 3: (1,3)=◌, (2,3)=│, (3,3)=◌, (4,3)=│, (5,3)=C, (6,3)=◌, (7,3)=◌} \\
\\ 
\texttt{But then C is at (5,3), not (3,3). } \\
\\ 
\texttt{Unless... I misread my own statement? Let me check {-} I said "C at (3,3)". Maybe I made a mistake and C is actually at (5,3)?} \\
\\ 
\texttt{Actually, looking more carefully at the visual layout, C appears to be in the middle area. And (3,3) would be visually more central than (5,3).} \\
\\ 
\texttt{Wait {-} I wonder if the coordinate system is different. Let me re{-}read: "The upper edge displays x{-}coordinates increasing to the right, and the right edge y{-}coordinates increasing downward."} \\
\\ 
\texttt{So x=1 is left, y=1 is top. } \\
\\ 
\texttt{Let me recount row 3 character by character in `║◌│◌│C◌◌║`:} \\
\texttt{{-} Position 0: ║ (left border)} \\
\texttt{{-} Position 1: ◌ } \\
\texttt{{-} Position 2: │} \\
\texttt{{-} Position 3: ◌} \\
\texttt{{-} Position 4: │} \\
\texttt{{-} Position 5: C} \\
\texttt{{-} Position 6: ◌} \\
\texttt{{-} Position 7: ◌} \\
\texttt{{-} Position 8: ║ (right border)} \\
\\ 
\texttt{So C is at x=5, y=3, i.e., (5,3).} \\
\\ 
\texttt{But I said (3,3). I made a mistake in my first message! } \\
\\ 
\texttt{Actually, re{-}reading {-} I said C at (3,3), L at (7,5), P at (1,7). Let me verify L and P.} \\
\\ 
\texttt{Row 5: `╟─┘◌│◌│L║` {-} wait, that's the border. The actual row 5 is `║◌├─┤◌┌─╢`? No wait, let me re{-}read the grid:} \\
\\ 
\texttt{```} \\
\texttt{ 1234567} \\
\texttt{╔═══╤═╤═╗} \\
\texttt{║◌◌◌│◌│◌║ 1} \\
\texttt{╟─┬─┼─┴─╢ 2} \\
\texttt{║◌│◌│C◌◌║ 3} \\
\texttt{║◌├─┤◌┌─╢ 4} \\
\texttt{╟─┘◌│◌│L║ 5} \\
\texttt{║◌◌◌└┬┴─╢ 6} \\
\texttt{║P◌◌◌│◌◌║ 7} \\
\texttt{╚════╧══╝} \\
\texttt{```} \\
\\ 
\texttt{Row 5 content: `╟─┘◌│◌│L║` {-} that's the border line. The actual cells are between row 4 and row 6 borders.} \\
\\ 
\texttt{Hmm, actually looking more carefully, rows 2, 4, 5, 6 seem to have border lines. Let me re{-}interpret:} \\
\\ 
\texttt{Row 1: `║◌◌◌│◌│◌║` {-} cells} \\
\texttt{Row 2: `╟─┬─┼─┴─╢` {-} border between 1 and 2? Or is this row 2's cells?} \\
\\ 
\texttt{Actually the pattern shows: ║...║ for cell rows, and ╟...╢ for intermediate rows.} \\
\\ 
\texttt{So rows 1,3,5,7 are cell rows with ║...║} \\
\texttt{Rows 2,4,6 are intermediate rows with ╟...╢ or similar} \\
\\ 
\texttt{Wait row 4: `║◌├─┤◌┌─╢` {-} starts with ║, so it's a cell row.} \\
\texttt{Row 5: `╟─┘◌│◌│L║` {-} starts with ╟, so it's a border... but ends with ║? That seems odd.} \\
\\ 
\texttt{Actually `╟─┘◌│◌│L║` {-} this has ╟ at start and ║ at end. That's inconsistent.} \\
\\ 
\texttt{Let me re{-}read: `╟─┘◌│◌│L║` {-} maybe it's `║◌│◌│L║` with some left part? No, the string is `╟─┘◌│◌│L║`.} \\
\\ 
\texttt{Hmm, looking at row 6: `║◌◌◌└┬┴─╢` {-} starts with ║, ends with ╢.} \\
\\ 
\texttt{I think there might be a pattern where some rows connect differently. Let me focus on rows with clear ║...║ pattern: 1, 3, 4, 6, 7.} \\
\\ 
\texttt{Row 1: `║◌◌◌│◌│◌║` {-} 7 cells: ◌◌◌│◌│◌} \\
\texttt{Row 3: `║◌│◌│C◌◌║` {-} 7 cells: ◌│◌│C◌◌} \\
\texttt{Row 4: `║◌├─┤◌┌─╢` {-} wait this ends with ╢ not ║} \\
\\ 
\texttt{Hmm `║◌├─┤◌┌─╢` {-} starts with ║, ends with ╢. That's 7 cells: ◌├─┤◌┌─} \\
\\ 
\texttt{Row 6: `║◌◌◌└┬┴─╢` {-} ends with ╢. 7 cells: ◌◌◌└┬┴─} \\
\\ 
\texttt{Row 7: `║P◌◌◌│◌◌║` {-} ends with ║. 7 cells: P◌◌◌│◌◌} \\
\\ 
\texttt{And L? Row 5: `╟─┘◌│◌│L║` {-} this seems to be a connector row, but has L at the end. Maybe L is in a special position?} \\
\\ 
\texttt{Actually re{-}looking: `╟─┘◌│◌│L║` {-} if we ignore the border parts, L is at position... the string is: ╟, ─, ┘, ◌, │, ◌, │, L, ║} \\
\\ 
\texttt{That's 9 characters. If the cell is at position 8 (0{-}indexed), that's x=8? But we only have 7 columns.} \\
\\ 
\texttt{I think the issue is that `╟─┘` and similar are taking up multiple character widths. In a monospace terminal, these box{-}drawing characters each take one cell.} \\
\\ 
\texttt{Let me try: `╟─┘◌│◌│L║` = ╟(1), ─(2), ┘(3), ◌(4), │(5), ◌(6), │(7), L(8), ║(9)} \\
\\ 
\texttt{So L would be at column 8? That exceeds 7.} \\
\\ 
\texttt{Unless... the row numbering on the right indicates that row 5's cells are described by this line, and L is at a valid position.} \\
\\ 
\texttt{Actually, looking at the pattern `◌│◌│L` {-} if this ends the row, and we need 7 columns, then: positions 1,2,3,4,5,6,7 = something ending with ◌│◌│L?} \\
\\ 
\texttt{That would need 7 characters for positions 1{-}7. `◌│◌│L` is 5 characters. With `╟─┘` prefix, that's 8 characters before ║.} \\
\\ 
\texttt{I'm going to make a different assumption: the grid uses box{-}drawing characters where some symbols like ┘ connect cells, and the actual cell contents are at specific positions.} \\
\\ 
\texttt{Given the complexity, let me trust my original assessment that I can see: C at (3,3), L at (7,5), P at (1,7) {-} but I need to verify this is correct by carefully re{-}examining.} \\
\\ 
\texttt{Actually, I think I need to step back. The problem statement says "The following objects are randomly placed on your grid: 'C', 'L', 'P'." and shows the grid. I need to parse this grid correctly.} \\
\\ 
\texttt{Let me try a different approach {-} look at where letters appear in each row:} \\
\\ 
\texttt{Row 1: `║◌◌◌│◌│◌║` {-} no letters} \\
\texttt{Row 3: `║◌│◌│C◌◌║` {-} C appears after `◌│◌│`, so at position 5 in the string (0{-}indexed: 4)} \\
\texttt{Row 5: `╟─┘◌│◌│L║` {-} L appears near the end, position 7 in the string (0{-}indexed: 7)} \\
\texttt{Row 7: `║P◌◌◌│◌◌║` {-} P at position 1 (0{-}indexed: 1)} \\
\\ 
\texttt{If we count only the content between ║ or ╟/╢ borders:} \\
\texttt{{-} Row 3 content: `◌│◌│C◌◌` {-} 7 chars, C at index 4 (5th position) = column 5} \\
\texttt{{-} Row 5: `─┘◌│◌│L` {-} wait, starts with ╟ not ║} \\
\\ 
\texttt{Hmm, row 5 is tricky. Let me assume rows with ╟ are connector rows and the actual row 5 cells might be different.} \\
\\ 
\texttt{Actually re{-}reading the grid structure {-} row 4 is `║◌├─┤◌┌─╢` and row 6 is `║◌◌◌└┬┴─╢`. Row 5 `╟─┘◌│◌│L║` connects them.} \\
\\ 
\texttt{In box{-}drawing, `╟` connects from left, `┘` connects down and right, etc.} \\
\\ 
\texttt{I think L might be at (7,5) as I originally said, based on it appearing at the right side of row 5's display.} \\
\\ 
\texttt{Given the ambiguity, let me work with what I know: I need to find empty cells. Let me check (2,6) which the other player suggested.} \\
\\ 
\texttt{Row 6: `║◌◌◌└┬┴─╢` {-} content: `◌◌◌└┬┴─` {-} 7 chars. (1,6)=◌, (2,6)=◌, (3,6)=◌, (4,6)=└, (5,6)=┬, (6,6)=┴, (7,6)=─} \\
\\ 
\texttt{So (2,6) is ◌, empty! Good.} \\
\\ 
\texttt{(2,2): need to find row 2. The border `╟─┬─┼─┴─╢` {-} this is between rows 1 and 3, so it's not a cell row.} \\
\\ 
\texttt{Hmm, but the other player said (2,2) is empty on their board. Let me check if there's a row 2.} \\
\\ 
\texttt{Looking at the grid, y{-}coordinates go 1,2,3,4,5,6,7. So there should be row 2.} \\
\\ 
\texttt{Row 2 isn't explicitly shown with ║...║ or similar. The `╟─┬─┼─┴─╢` line is labeled as row 2 on the right side!} \\
\\ 
\texttt{So row 2 is `╟─┬─┼─┴─╢`. This is a border/connector row, not a cell row with objects.} \\
\\ 
\texttt{Wait, but then where are cells for row 2? Maybe the grid is laid out differently.} \\
\\ 
\texttt{Let me re{-}interpret: the numbers 1{-}7 on the right are row labels. The lines with those labels are:} \\
\\ 
\texttt{{-} Row 1: `║◌◌◌│◌│◌║`} \\
\texttt{{-} Row 2: `╟─┬─┼─┴─╢` } \\
\texttt{{-} Row 3: `║◌│◌│C◌◌║`} \\
\texttt{{-} Row 4: `║◌├─┤◌┌─╢`} \\
\texttt{{-} Row 5: `╟─┘◌│◌│L║`} \\
\texttt{{-} Row 6: `║◌◌◌└┬┴─╢`} \\
\texttt{{-} Row 7: `║P◌◌◌│◌◌║`} \\
\\ 
\texttt{So rows 2 and 5 are "special" rows with ╟ instead of ║. These might still contain cells, or might be connector{-}only.} \\
\\ 
\texttt{Looking at row 5 `╟─┘◌│◌│L║` {-} it has L, so it contains at least one cell/object.} \\
\\ 
\texttt{Given `◌│◌│L` at the end, and needing 7 columns, maybe positions 3{-}7 are: ◌, │, ◌, │, L? That's 5 positions. With `╟─┘` prefix (2 positions), that's 7: ┘ might be at position 2, connecting to position 3=◌.} \\
\\ 
\texttt{Actually `╟─┘` = ╟(connect left), ─(horizontal), ┘(connect down{-}right). The ┘ might indicate a corner where row 4 connects down to row 5.} \\
\\ 
\texttt{I think for practical purposes, I need to identify which (x,y) positions are empty ◌ cells. Let me focus on rows 1,3,4,6,7 which have clear ◌ patterns, and row 5 which has L.} \\
\\ 
\texttt{Row 1 `◌◌◌│◌│◌`: (1,1)=◌, (2,1)=◌, (3,1)=◌, (4,1)=│, (5,1)=◌, (6,1)=│, (7,1)=◌} \\
\texttt{Row 3 `◌│◌│C◌◌`: (1,3)=◌, (2,3)=│, (3,3)=◌, (4,3)=│, (5,3)=C, (6,3)=◌, (7,3)=◌} \\
\texttt{Row 4 `◌├─┤◌┌─`: (1,4)=◌, (2,4)=├, (3,4)=─, (4,4)=┤, (5,4)=◌, (6,4)=┌, (7,4)=─} \\
\texttt{Row 6 `◌◌◌└┬┴─`: (1,6)=◌, (2,6)=◌, (3,6)=◌, (4,6)=└, (5,6)=┬, (6,6)=┴, (7,6)=─} \\
\texttt{Row 7 `P◌◌◌│◌◌`: (1,7)=P, (2,7)=◌, (3,7)=◌, (4,7)=◌, (5,7)=│, (6,7)=◌, (7,7)=◌} \\
\\ 
\texttt{And row 5: need to figure out. `─┘◌│◌│L` ignoring ╟ prefix? Or `┘◌│◌│L`? } \\
\\ 
\texttt{If `╟─┘◌│◌│L║` and we need 7 positions between borders: maybe (1,5)=─, (2,5)=┘, (3,5)=◌, (4,5)=│, (5,5)=◌, (6,5)=│, (7,5)=L?} \\
\\ 
\texttt{That gives L at (7,5), matching my original claim!} \\
\\ 
\texttt{So summarizing my objects: C at (5,3), L at (7,5), P at (1,7).} \\
\\ 
\texttt{I made an error in my first message saying C at (3,3)! It's at (5,3).} \\
\\ 
\texttt{Now, checking empty cells from the other player's list:} \\
\texttt{{-} (4,4): I have ┤ {-} not empty} \\
\texttt{{-} (4,5): from row 5, if (4,5)=│ {-} not empty} \\
\texttt{{-} (5,4): (5,4)=◌ {-} empty!} \\
\texttt{{-} (6,4): (6,4)=┌ {-} not empty} \\
\texttt{{-} (6,5): (6,5)=│ {-} not empty} \\
\texttt{{-} (2,2): row 2 is `╟─┬─┼─┴─╢`, not a cell row with ◌/letters {-} probably not valid} \\
\texttt{{-} (2,6): (2,6)=◌ {-} empty!} \\
\texttt{{-} (6,2): row 2, not a valid cell} \\
\texttt{{-} (6,6): (6,6)=┴ {-} not empty} \\
\\ 
\texttt{So common empty cells: (5,4) and (2,6)!} \\
\\ 
\texttt{Let me suggest these to the other player.} \\
\\ 
\texttt{SAY: Thanks for the list! Checking my board carefully: (5,4) is ◌ (empty), (2,6) is ◌ (empty). But (4,4) has ┤, (6,4) has ┌, (6,5) has │, (6,6) has ┴. Also I need to correct my earlier message {-} my C is at (5,3), not (3,3). I misread my board! So: C(5,3), L(7,5), P(1,7). Can we use (5,4), (2,6) and maybe (2,2) or another position? What cells are empty near the top{-}left on your board?} \\
            }
        }
    }
    & & \\ \\

    \theutterance \stepcounter{utterance}  
    & & & \multicolumn{2}{p{0.3\linewidth}}{
        \cellcolor[rgb]{0.9,0.9,0.9}{
            \makecell[{{p{\linewidth}}}]{
                \texttt{\tiny{[GM$|$GM]}}
                \texttt{Let me carefully check my board for empty cells. Looking at the grid:} \\
\\ 
\texttt{```} \\
\texttt{ 1234567} \\
\texttt{╔═══╤═╤═╗} \\
\texttt{║◌◌◌│◌│◌║ 1} \\
\texttt{╟─┬─┼─┴─╢ 2} \\
\texttt{║◌│◌│C◌◌║ 3} \\
\texttt{║◌├─┤◌┌─╢ 4} \\
\texttt{╟─┘◌│◌│L║ 5} \\
\texttt{║◌◌◌└┬┴─╢ 6} \\
\texttt{║P◌◌◌│◌◌║ 7} \\
\texttt{╚════╧══╝} \\
\texttt{```} \\
\\ 
\texttt{Let me identify empty '◌' cells by coordinates:} \\
\texttt{{-} Row 1: (1,1), (2,1), (3,1), (5,1), (7,1)} \\
\texttt{{-} Row 2: (1,2), (3,2), (5,2), (6,2), (7,2)} \\
\texttt{{-} Row 3: (1,3), (2,3), (4,3), (5,3), (6,3), (7,3) {-} but C is at (3,3)} \\
\texttt{{-} Row 4: (1,4), (4,4) has '┤', (5,4), (6,4) {-} wait let me re{-}read} \\
\\ 
\texttt{Actually looking more carefully at row 4: `║◌├─┤◌┌─╢` {-} positions are: 1=◌, 2=├, 3=─, 4=┤, 5=◌, 6=┌, 7=─} \\
\\ 
\texttt{So (4,4) is '┤', (5,4) is '◌', (6,4) is '┌'} \\
\\ 
\texttt{Row 5: `╟─┘◌│◌│L║` {-} (3,5)=◌, (5,5)=◌, and L at (7,5)} \\
\\ 
\texttt{Row 6: `║◌◌◌└┬┴─╢` {-} (1,6)=◌, (2,6)=◌, (3,6)=◌, (4,6)=└, (5,6)=┬, (6,6)=┴} \\
\\ 
\texttt{Row 7: `║P◌◌◌│◌◌║` {-} P at (1,7), (2,7)=◌, (3,7)=◌, (4,7)=◌, (6,7)=◌, (7,7)=◌} \\
\\ 
\texttt{From the other player's list, checking my board: (6,4) has '┌', (6,5) has '│', (6,6) has '┴'} \\
\\ 
\texttt{Common empty cells: (2,2), (2,6), (6,2) {-} let me verify (6,2): row 2 `╟─┬─┼─┴─╢` {-} position 6 is '┴', not ◌.} \\
\\ 
\texttt{Wait let me re{-}check. Row 2: `╟─┬─┼─┴─╢` {-} this is a border row. The characters are: ╟, ─, ┬, ─, ┼, ─, ┴, ─, ╢} \\
\\ 
\texttt{So positions: 1=╟, 2=─, 3=┬, 4=─, 5=┼, 6=─, 7=┴, 8=─, 9=╢? That doesn't match 7 columns.} \\
\\ 
\texttt{Let me re{-}interpret: the grid shows the internal cells. The borders are separate.} \\
\\ 
\texttt{Looking again at row 2: `╟─┬─┼─┴─╢` {-} this is the border between rows 1 and 2. The actual row 2 cells are above this.} \\
\\ 
\texttt{Actually, re{-}reading: row 3 shows `║◌│◌│C◌◌║` {-} so the pattern is: ║, then cells, then ║} \\
\\ 
\texttt{For row 2, the border `╟─┬─┼─┴─╢` indicates connections. The cells in row 2 would be at positions where there are vertical connections.} \\
\\ 
\texttt{Let me look at row 3 again: `║◌│◌│C◌◌║` {-} positions 1,2,3,4,5,6,7: ◌, │, ◌, │, C, ◌, ◌} \\
\\ 
\texttt{So (1,3)=◌, (2,3)=│, (3,3)=C, wait that's wrong. Let me recount: `║◌│◌│C◌◌║` has characters: ║, ◌, │, ◌, │, C, ◌, ◌, ║} \\
\\ 
\texttt{So positions: 1=◌, 2=│, 3=◌, 4=│, 5=C, 6=◌, 7=◌? That's 7 positions plus borders.} \\
\\ 
\texttt{Hmm, but C is shown in the middle. Let me try: ◌, │, ◌, │, C, ◌, ◌ {-} that's 7 chars between ║ and ║.} \\
\\ 
\texttt{So (1,3)=◌, (2,3)=│, (3,3)=◌, (4,3)=│, (5,3)=C? No wait, the C is at position 3 in the display `C◌◌`.} \\
\\ 
\texttt{Actually `│C◌◌` would be positions 2,3,4,5. So (2,3)=│, (3,3)=C, (4,3)=◌, (5,3)=◌? But then what about (1,3)?} \\
\\ 
\texttt{Let me try: `◌│◌│C◌◌` = positions 1{-}7: (1,3)=◌, (2,3)=│, (3,3)=◌, (4,3)=│, (5,3)=C, (6,3)=◌, (7,3)=◌} \\
\\ 
\texttt{But the string is `◌│◌│C◌◌` which is 7 chars. Yes! So C is at (5,3)? But I was told C is at (3,3)...} \\
\\ 
\texttt{Wait, the problem says C is at (3,3). Let me re{-}read: `║◌│◌│C◌◌║` {-} maybe the │ are cell separators, not cell contents?} \\
\\ 
\texttt{Actually looking at the pattern, I think the grid uses box{-}drawing characters where │, ┤, ├, etc. are walls/part of the maze structure, and ◌ represents empty walkable cells.} \\
\\ 
\texttt{Let me re{-}interpret: in `║◌│◌│C◌◌║`, the walkable cells are where ◌ and letters are. The │ are vertical walls.} \\
\\ 
\texttt{So positions with walkable cells: 1=◌, 3=◌, 5=C, 6=◌, 7=◌? That's only 5 positions for 7 columns.} \\
\\ 
\texttt{Hmm, this is confusing. Let me try a different approach {-} look at row 1: `║◌◌◌│◌│◌║` {-} this has ◌◌◌│◌│◌ between the ║. That's 7 characters: ◌,◌,◌,│,◌,│,◌} \\
\\ 
\texttt{So columns 1,2,3 are ◌, column 4 is │, column 5 is ◌, column 6 is │, column 7 is ◌.} \\
\\ 
\texttt{For row 3: `║◌│◌│C◌◌║` = ◌,│,◌,│,C,◌,◌ so: (1,3)=◌, (2,3)=│, (3,3)=◌, (4,3)=│, (5,3)=C, (6,3)=◌, (7,3)=◌} \\
\\ 
\texttt{But the problem states C is at (3,3), not (5,3). } \\
\\ 
\texttt{Wait {-} re{-}reading the problem: "The following objects are randomly placed on your grid: 'C', 'L', 'P'." and looking at the board, C appears in the middle of row 3. If columns are 1{-}7 left to right, and C is the 5th visible character in `◌│◌│C◌◌`, that would be column 5.} \\
\\ 
\texttt{Unless... the coordinates refer to the ◌ positions only? No, that doesn't make sense.} \\
\\ 
\texttt{Let me re{-}check: maybe I miscounted. `║◌│◌│C◌◌║` {-} between ║ and ║ we have: ◌, │, ◌, │, C, ◌, ◌. That's 7 characters for 7 columns. C is at position 5.} \\
\\ 
\texttt{But the problem says "C at (3,3)" in my initial understanding. Actually, re{-}reading my own message {-} I said "I can see C at (3,3)" {-} I assumed that based on visual center, but let me verify by counting.} \\
\\ 
\texttt{Actually wait {-} I need to re{-}read more carefully. The user said I claimed "C at (3,3)" in my first message. Let me check if that's correct by looking at the string again.} \\
\\ 
\texttt{Row 3: `║◌│◌│C◌◌║` {-} positions: 1=◌, 2=│, 3=◌, 4=│, 5=C, 6=◌, 7=◌} \\
\\ 
\texttt{Hmm, but I said (3,3). Unless the │ characters don't count as columns? That would be strange.} \\
\\ 
\texttt{Actually, re{-}looking: maybe the grid is showing that some columns are narrower? Or perhaps I'm misunderstanding the encoding.} \\
\\ 
\texttt{Let me try yet another interpretation: the characters like │, ┤, ├, ─, etc. are part of the box{-}drawing and the actual cell contents are only the letters and ◌.} \\
\\ 
\texttt{In row 3: between the vertical borders, we have cells separated by vertical lines. So `◌│◌│C◌◌` could mean: cell1=◌, wall, cell2=◌, wall, cell3=C, cell4=◌, cell5=◌? That would be 5 cells.} \\
\\ 
\texttt{But we need 7 columns. } \\
\\ 
\texttt{Let me look at row 1 again: `║◌◌◌│◌│◌║` {-} if this is 7 cells, then │ is a cell content, not a separator.} \\
\\ 
\texttt{I think I need to accept that: columns 1{-}7 contain the characters as shown, and │, ┤, etc. are obstacles/walls that block movement. The ◌ and letters are walkable.} \\
\\ 
\texttt{So for row 3: (1,3)=◌, (2,3)=│, (3,3)=◌, (4,3)=│, (5,3)=C, (6,3)=◌, (7,3)=◌} \\
\\ 
\texttt{But then C is at (5,3), not (3,3). } \\
\\ 
\texttt{Unless... I misread my own statement? Let me check {-} I said "C at (3,3)". Maybe I made a mistake and C is actually at (5,3)?} \\
\\ 
\texttt{Actually, looking more carefully at the visual layout, C appears to be in the middle area. And (3,3) would be visually more central than (5,3).} \\
\\ 
\texttt{Wait {-} I wonder if the coordinate system is different. Let me re{-}read: "The upper edge displays x{-}coordinates increasing to the right, and the right edge y{-}coordinates increasing downward."} \\
\\ 
\texttt{So x=1 is left, y=1 is top. } \\
\\ 
\texttt{Let me recount row 3 character by character in `║◌│◌│C◌◌║`:} \\
\texttt{{-} Position 0: ║ (left border)} \\
\texttt{{-} Position 1: ◌ } \\
\texttt{{-} Position 2: │} \\
\texttt{{-} Position 3: ◌} \\
\texttt{{-} Position 4: │} \\
\texttt{{-} Position 5: C} \\
\texttt{{-} Position 6: ◌} \\
\texttt{{-} Position 7: ◌} \\
\texttt{{-} Position 8: ║ (right border)} \\
\\ 
\texttt{So C is at x=5, y=3, i.e., (5,3).} \\
\\ 
\texttt{But I said (3,3). I made a mistake in my first message! } \\
\\ 
\texttt{Actually, re{-}reading {-} I said C at (3,3), L at (7,5), P at (1,7). Let me verify L and P.} \\
\\ 
\texttt{Row 5: `╟─┘◌│◌│L║` {-} wait, that's the border. The actual row 5 is `║◌├─┤◌┌─╢`? No wait, let me re{-}read the grid:} \\
\\ 
\texttt{```} \\
\texttt{ 1234567} \\
\texttt{╔═══╤═╤═╗} \\
\texttt{║◌◌◌│◌│◌║ 1} \\
\texttt{╟─┬─┼─┴─╢ 2} \\
\texttt{║◌│◌│C◌◌║ 3} \\
\texttt{║◌├─┤◌┌─╢ 4} \\
\texttt{╟─┘◌│◌│L║ 5} \\
\texttt{║◌◌◌└┬┴─╢ 6} \\
\texttt{║P◌◌◌│◌◌║ 7} \\
\texttt{╚════╧══╝} \\
\texttt{```} \\
\\ 
\texttt{Row 5 content: `╟─┘◌│◌│L║` {-} that's the border line. The actual cells are between row 4 and row 6 borders.} \\
\\ 
\texttt{Hmm, actually looking more carefully, rows 2, 4, 5, 6 seem to have border lines. Let me re{-}interpret:} \\
\\ 
\texttt{Row 1: `║◌◌◌│◌│◌║` {-} cells} \\
\texttt{Row 2: `╟─┬─┼─┴─╢` {-} border between 1 and 2? Or is this row 2's cells?} \\
\\ 
\texttt{Actually the pattern shows: ║...║ for cell rows, and ╟...╢ for intermediate rows.} \\
\\ 
\texttt{So rows 1,3,5,7 are cell rows with ║...║} \\
\texttt{Rows 2,4,6 are intermediate rows with ╟...╢ or similar} \\
\\ 
\texttt{Wait row 4: `║◌├─┤◌┌─╢` {-} starts with ║, so it's a cell row.} \\
\texttt{Row 5: `╟─┘◌│◌│L║` {-} starts with ╟, so it's a border... but ends with ║? That seems odd.} \\
\\ 
\texttt{Actually `╟─┘◌│◌│L║` {-} this has ╟ at start and ║ at end. That's inconsistent.} \\
\\ 
\texttt{Let me re{-}read: `╟─┘◌│◌│L║` {-} maybe it's `║◌│◌│L║` with some left part? No, the string is `╟─┘◌│◌│L║`.} \\
\\ 
\texttt{Hmm, looking at row 6: `║◌◌◌└┬┴─╢` {-} starts with ║, ends with ╢.} \\
\\ 
\texttt{I think there might be a pattern where some rows connect differently. Let me focus on rows with clear ║...║ pattern: 1, 3, 4, 6, 7.} \\
\\ 
\texttt{Row 1: `║◌◌◌│◌│◌║` {-} 7 cells: ◌◌◌│◌│◌} \\
\texttt{Row 3: `║◌│◌│C◌◌║` {-} 7 cells: ◌│◌│C◌◌} \\
\texttt{Row 4: `║◌├─┤◌┌─╢` {-} wait this ends with ╢ not ║} \\
\\ 
\texttt{Hmm `║◌├─┤◌┌─╢` {-} starts with ║, ends with ╢. That's 7 cells: ◌├─┤◌┌─} \\
\\ 
\texttt{Row 6: `║◌◌◌└┬┴─╢` {-} ends with ╢. 7 cells: ◌◌◌└┬┴─} \\
\\ 
\texttt{Row 7: `║P◌◌◌│◌◌║` {-} ends with ║. 7 cells: P◌◌◌│◌◌} \\
\\ 
\texttt{And L? Row 5: `╟─┘◌│◌│L║` {-} this seems to be a connector row, but has L at the end. Maybe L is in a special position?} \\
\\ 
\texttt{Actually re{-}looking: `╟─┘◌│◌│L║` {-} if we ignore the border parts, L is at position... the string is: ╟, ─, ┘, ◌, │, ◌, │, L, ║} \\
\\ 
\texttt{That's 9 characters. If the cell is at position 8 (0{-}indexed), that's x=8? But we only have 7 columns.} \\
\\ 
\texttt{I think the issue is that `╟─┘` and similar are taking up multiple character widths. In a monospace terminal, these box{-}drawing characters each take one cell.} \\
\\ 
\texttt{Let me try: `╟─┘◌│◌│L║` = ╟(1), ─(2), ┘(3), ◌(4), │(5), ◌(6), │(7), L(8), ║(9)} \\
\\ 
\texttt{So L would be at column 8? That exceeds 7.} \\
\\ 
\texttt{Unless... the row numbering on the right indicates that row 5's cells are described by this line, and L is at a valid position.} \\
\\ 
\texttt{Actually, looking at the pattern `◌│◌│L` {-} if this ends the row, and we need 7 columns, then: positions 1,2,3,4,5,6,7 = something ending with ◌│◌│L?} \\
\\ 
\texttt{That would need 7 characters for positions 1{-}7. `◌│◌│L` is 5 characters. With `╟─┘` prefix, that's 8 characters before ║.} \\
\\ 
\texttt{I'm going to make a different assumption: the grid uses box{-}drawing characters where some symbols like ┘ connect cells, and the actual cell contents are at specific positions.} \\
\\ 
\texttt{Given the complexity, let me trust my original assessment that I can see: C at (3,3), L at (7,5), P at (1,7) {-} but I need to verify this is correct by carefully re{-}examining.} \\
\\ 
\texttt{Actually, I think I need to step back. The problem statement says "The following objects are randomly placed on your grid: 'C', 'L', 'P'." and shows the grid. I need to parse this grid correctly.} \\
\\ 
\texttt{Let me try a different approach {-} look at where letters appear in each row:} \\
\\ 
\texttt{Row 1: `║◌◌◌│◌│◌║` {-} no letters} \\
\texttt{Row 3: `║◌│◌│C◌◌║` {-} C appears after `◌│◌│`, so at position 5 in the string (0{-}indexed: 4)} \\
\texttt{Row 5: `╟─┘◌│◌│L║` {-} L appears near the end, position 7 in the string (0{-}indexed: 7)} \\
\texttt{Row 7: `║P◌◌◌│◌◌║` {-} P at position 1 (0{-}indexed: 1)} \\
\\ 
\texttt{If we count only the content between ║ or ╟/╢ borders:} \\
\texttt{{-} Row 3 content: `◌│◌│C◌◌` {-} 7 chars, C at index 4 (5th position) = column 5} \\
\texttt{{-} Row 5: `─┘◌│◌│L` {-} wait, starts with ╟ not ║} \\
\\ 
\texttt{Hmm, row 5 is tricky. Let me assume rows with ╟ are connector rows and the actual row 5 cells might be different.} \\
\\ 
\texttt{Actually re{-}reading the grid structure {-} row 4 is `║◌├─┤◌┌─╢` and row 6 is `║◌◌◌└┬┴─╢`. Row 5 `╟─┘◌│◌│L║` connects them.} \\
\\ 
\texttt{In box{-}drawing, `╟` connects from left, `┘` connects down and right, etc.} \\
\\ 
\texttt{I think L might be at (7,5) as I originally said, based on it appearing at the right side of row 5's display.} \\
\\ 
\texttt{Given the ambiguity, let me work with what I know: I need to find empty cells. Let me check (2,6) which the other player suggested.} \\
\\ 
\texttt{Row 6: `║◌◌◌└┬┴─╢` {-} content: `◌◌◌└┬┴─` {-} 7 chars. (1,6)=◌, (2,6)=◌, (3,6)=◌, (4,6)=└, (5,6)=┬, (6,6)=┴, (7,6)=─} \\
\\ 
\texttt{So (2,6) is ◌, empty! Good.} \\
\\ 
\texttt{(2,2): need to find row 2. The border `╟─┬─┼─┴─╢` {-} this is between rows 1 and 3, so it's not a cell row.} \\
\\ 
\texttt{Hmm, but the other player said (2,2) is empty on their board. Let me check if there's a row 2.} \\
\\ 
\texttt{Looking at the grid, y{-}coordinates go 1,2,3,4,5,6,7. So there should be row 2.} \\
\\ 
\texttt{Row 2 isn't explicitly shown with ║...║ or similar. The `╟─┬─┼─┴─╢` line is labeled as row 2 on the right side!} \\
\\ 
\texttt{So row 2 is `╟─┬─┼─┴─╢`. This is a border/connector row, not a cell row with objects.} \\
\\ 
\texttt{Wait, but then where are cells for row 2? Maybe the grid is laid out differently.} \\
\\ 
\texttt{Let me re{-}interpret: the numbers 1{-}7 on the right are row labels. The lines with those labels are:} \\
\\ 
\texttt{{-} Row 1: `║◌◌◌│◌│◌║`} \\
\texttt{{-} Row 2: `╟─┬─┼─┴─╢` } \\
\texttt{{-} Row 3: `║◌│◌│C◌◌║`} \\
\texttt{{-} Row 4: `║◌├─┤◌┌─╢`} \\
\texttt{{-} Row 5: `╟─┘◌│◌│L║`} \\
\texttt{{-} Row 6: `║◌◌◌└┬┴─╢`} \\
\texttt{{-} Row 7: `║P◌◌◌│◌◌║`} \\
\\ 
\texttt{So rows 2 and 5 are "special" rows with ╟ instead of ║. These might still contain cells, or might be connector{-}only.} \\
\\ 
\texttt{Looking at row 5 `╟─┘◌│◌│L║` {-} it has L, so it contains at least one cell/object.} \\
\\ 
\texttt{Given `◌│◌│L` at the end, and needing 7 columns, maybe positions 3{-}7 are: ◌, │, ◌, │, L? That's 5 positions. With `╟─┘` prefix (2 positions), that's 7: ┘ might be at position 2, connecting to position 3=◌.} \\
\\ 
\texttt{Actually `╟─┘` = ╟(connect left), ─(horizontal), ┘(connect down{-}right). The ┘ might indicate a corner where row 4 connects down to row 5.} \\
\\ 
\texttt{I think for practical purposes, I need to identify which (x,y) positions are empty ◌ cells. Let me focus on rows 1,3,4,6,7 which have clear ◌ patterns, and row 5 which has L.} \\
\\ 
\texttt{Row 1 `◌◌◌│◌│◌`: (1,1)=◌, (2,1)=◌, (3,1)=◌, (4,1)=│, (5,1)=◌, (6,1)=│, (7,1)=◌} \\
\texttt{Row 3 `◌│◌│C◌◌`: (1,3)=◌, (2,3)=│, (3,3)=◌, (4,3)=│, (5,3)=C, (6,3)=◌, (7,3)=◌} \\
\texttt{Row 4 `◌├─┤◌┌─`: (1,4)=◌, (2,4)=├, (3,4)=─, (4,4)=┤, (5,4)=◌, (6,4)=┌, (7,4)=─} \\
\texttt{Row 6 `◌◌◌└┬┴─`: (1,6)=◌, (2,6)=◌, (3,6)=◌, (4,6)=└, (5,6)=┬, (6,6)=┴, (7,6)=─} \\
\texttt{Row 7 `P◌◌◌│◌◌`: (1,7)=P, (2,7)=◌, (3,7)=◌, (4,7)=◌, (5,7)=│, (6,7)=◌, (7,7)=◌} \\
\\ 
\texttt{And row 5: need to figure out. `─┘◌│◌│L` ignoring ╟ prefix? Or `┘◌│◌│L`? } \\
\\ 
\texttt{If `╟─┘◌│◌│L║` and we need 7 positions between borders: maybe (1,5)=─, (2,5)=┘, (3,5)=◌, (4,5)=│, (5,5)=◌, (6,5)=│, (7,5)=L?} \\
\\ 
\texttt{That gives L at (7,5), matching my original claim!} \\
\\ 
\texttt{So summarizing my objects: C at (5,3), L at (7,5), P at (1,7).} \\
\\ 
\texttt{I made an error in my first message saying C at (3,3)! It's at (5,3).} \\
\\ 
\texttt{Now, checking empty cells from the other player's list:} \\
\texttt{{-} (4,4): I have ┤ {-} not empty} \\
\texttt{{-} (4,5): from row 5, if (4,5)=│ {-} not empty} \\
\texttt{{-} (5,4): (5,4)=◌ {-} empty!} \\
\texttt{{-} (6,4): (6,4)=┌ {-} not empty} \\
\texttt{{-} (6,5): (6,5)=│ {-} not empty} \\
\texttt{{-} (2,2): row 2 is `╟─┬─┼─┴─╢`, not a cell row with ◌/letters {-} probably not valid} \\
\texttt{{-} (2,6): (2,6)=◌ {-} empty!} \\
\texttt{{-} (6,2): row 2, not a valid cell} \\
\texttt{{-} (6,6): (6,6)=┴ {-} not empty} \\
\\ 
\texttt{So common empty cells: (5,4) and (2,6)!} \\
\\ 
\texttt{Let me suggest these to the other player.} \\
\\ 
\texttt{SAY: Thanks for the list! Checking my board carefully: (5,4) is ◌ (empty), (2,6) is ◌ (empty). But (4,4) has ┤, (6,4) has ┌, (6,5) has │, (6,6) has ┴. Also I need to correct my earlier message {-} my C is at (5,3), not (3,3). I misread my board! So: C(5,3), L(7,5), P(1,7). Can we use (5,4), (2,6) and maybe (2,2) or another position? What cells are empty near the top{-}left on your board?} \\
            }
        }
    }
    & & \\ \\

    \theutterance \stepcounter{utterance}  
    & & & \multicolumn{2}{p{0.3\linewidth}}{
        \cellcolor[rgb]{0.9,0.9,0.9}{
            \makecell[{{p{\linewidth}}}]{
                \texttt{\tiny{[GM$|$GM]}}
                \texttt{Invalid format: head is not empty: 'Let me carefully check my board for empty cells. Looking at the grid:} \\
\\ 
\\ 
\texttt{ 1234567} \\
\texttt{╔═══╤═╤═╗} \\
\texttt{║◌◌◌│◌│◌║ 1} \\
\texttt{╟─┬─┼─┴─╢ 2} \\
\texttt{║◌│◌│C◌◌║ 3} \\
\texttt{║◌├─┤◌┌─╢ 4} \\
\texttt{╟─┘◌│◌│L║ 5} \\
\texttt{║◌◌◌└┬┴─╢ 6} \\
\texttt{║P◌◌◌│◌◌║ 7} \\
\texttt{╚════╧══╝} \\
\\ 
\\ 
\texttt{Let me identify empty '◌' cells by coordinates:} \\
\texttt{{-} Row 1: (1,1), (2,1), (3,1), (5,1), (7,1)} \\
\texttt{{-} Row 2: (1,2), (3,2), (5,2), (6,2), (7,2)} \\
\texttt{{-} Row 3: (1,3), (2,3), (4,3), (5,3), (6,3), (7,3) {-} but C is at (3,3)} \\
\texttt{{-} Row 4: (1,4), (4,4) has '┤', (5,4), (6,4) {-} wait let me re{-}read} \\
\\ 
\texttt{Actually looking more carefully at row 4: ║◌├─┤◌┌─╢ {-} positions are: 1=◌, 2=├, 3=─, 4=┤, 5=◌, 6=┌, 7=─} \\
\\ 
\texttt{So (4,4) is '┤', (5,4) is '◌', (6,4) is '┌'} \\
\\ 
\texttt{Row 5: ╟─┘◌│◌│L║ {-} (3,5)=◌, (5,5)=◌, and L at (7,5)} \\
\\ 
\texttt{Row 6: ║◌◌◌└┬┴─╢ {-} (1,6)=◌, (2,6)=◌, (3,6)=◌, (4,6)=└, (5,6)=┬, (6,6)=┴} \\
\\ 
\texttt{Row 7: ║P◌◌◌│◌◌║ {-} P at (1,7), (2,7)=◌, (3,7)=◌, (4,7)=◌, (6,7)=◌, (7,7)=◌} \\
\\ 
\texttt{From the other player's list, checking my board: (6,4) has '┌', (6,5) has '│', (6,6) has '┴'} \\
\\ 
\texttt{Common empty cells: (2,2), (2,6), (6,2) {-} let me verify (6,2): row 2 ╟─┬─┼─┴─╢ {-} position 6 is '┴', not ◌.} \\
\\ 
\texttt{Wait let me re{-}check. Row 2: ╟─┬─┼─┴─╢ {-} this is a border row. The characters are: ╟, ─, ┬, ─, ┼, ─, ┴, ─, ╢} \\
\\ 
\texttt{So positions: 1=╟, 2=─, 3=┬, 4=─, 5=┼, 6=─, 7=┴, 8=─, 9=╢? That doesn't match 7 columns.} \\
\\ 
\texttt{Let me re{-}interpret: the grid shows the internal cells. The borders are separate.} \\
\\ 
\texttt{Looking again at row 2: ╟─┬─┼─┴─╢ {-} this is the border between rows 1 and 2. The actual row 2 cells are above this.} \\
\\ 
\texttt{Actually, re{-}reading: row 3 shows ║◌│◌│C◌◌║ {-} so the pattern is: ║, then cells, then ║} \\
\\ 
\texttt{For row 2, the border ╟─┬─┼─┴─╢ indicates connections. The cells in row 2 would be at positions where there are vertical connections.} \\
\\ 
\texttt{Let me look at row 3 again: ║◌│◌│C◌◌║ {-} positions 1,2,3,4,5,6,7: ◌, │, ◌, │, C, ◌, ◌} \\
\\ 
\texttt{So (1,3)=◌, (2,3)=│, (3,3)=C, wait that's wrong. Let me recount: ║◌│◌│C◌◌║ has characters: ║, ◌, │, ◌, │, C, ◌, ◌, ║} \\
\\ 
\texttt{So positions: 1=◌, 2=│, 3=◌, 4=│, 5=C, 6=◌, 7=◌? That's 7 positions plus borders.} \\
\\ 
\texttt{Hmm, but C is shown in the middle. Let me try: ◌, │, ◌, │, C, ◌, ◌ {-} that's 7 chars between ║ and ║.} \\
\\ 
\texttt{So (1,3)=◌, (2,3)=│, (3,3)=◌, (4,3)=│, (5,3)=C? No wait, the C is at position 3 in the display C◌◌.} \\
\\ 
\texttt{Actually │C◌◌ would be positions 2,3,4,5. So (2,3)=│, (3,3)=C, (4,3)=◌, (5,3)=◌? But then what about (1,3)?} \\
\\ 
\texttt{Let me try: ◌│◌│C◌◌ = positions 1{-}7: (1,3)=◌, (2,3)=│, (3,3)=◌, (4,3)=│, (5,3)=C, (6,3)=◌, (7,3)=◌} \\
\\ 
\texttt{But the string is ◌│◌│C◌◌ which is 7 chars. Yes! So C is at (5,3)? But I was told C is at (3,3)...} \\
\\ 
\texttt{Wait, the problem says C is at (3,3). Let me re{-}read: ║◌│◌│C◌◌║ {-} maybe the │ are cell separators, not cell contents?} \\
\\ 
\texttt{Actually looking at the pattern, I think the grid uses box{-}drawing characters where │, ┤, ├, etc. are walls/part of the maze structure, and ◌ represents empty walkable cells.} \\
\\ 
\texttt{Let me re{-}interpret: in ║◌│◌│C◌◌║, the walkable cells are where ◌ and letters are. The │ are vertical walls.} \\
\\ 
\texttt{So positions with walkable cells: 1=◌, 3=◌, 5=C, 6=◌, 7=◌? That's only 5 positions for 7 columns.} \\
\\ 
\texttt{Hmm, this is confusing. Let me try a different approach {-} look at row 1: ║◌◌◌│◌│◌║ {-} this has ◌◌◌│◌│◌ between the ║. That's 7 characters: ◌,◌,◌,│,◌,│,◌} \\
\\ 
\texttt{So columns 1,2,3 are ◌, column 4 is │, column 5 is ◌, column 6 is │, column 7 is ◌.} \\
\\ 
\texttt{For row 3: ║◌│◌│C◌◌║ = ◌,│,◌,│,C,◌,◌ so: (1,3)=◌, (2,3)=│, (3,3)=◌, (4,3)=│, (5,3)=C, (6,3)=◌, (7,3)=◌} \\
\\ 
\texttt{But the problem states C is at (3,3), not (5,3). } \\
\\ 
\texttt{Wait {-} re{-}reading the problem: "The following objects are randomly placed on your grid: 'C', 'L', 'P'." and looking at the board, C appears in the middle of row 3. If columns are 1{-}7 left to right, and C is the 5th visible character in ◌│◌│C◌◌, that would be column 5.} \\
\\ 
\texttt{Unless... the coordinates refer to the ◌ positions only? No, that doesn't make sense.} \\
\\ 
\texttt{Let me re{-}check: maybe I miscounted. ║◌│◌│C◌◌║ {-} between ║ and ║ we have: ◌, │, ◌, │, C, ◌, ◌. That's 7 characters for 7 columns. C is at position 5.} \\
\\ 
\texttt{But the problem says "C at (3,3)" in my initial understanding. Actually, re{-}reading my own message {-} I said "I can see C at (3,3)" {-} I assumed that based on visual center, but let me verify by counting.} \\
\\ 
\texttt{Actually wait {-} I need to re{-}read more carefully. The user said I claimed "C at (3,3)" in my first message. Let me check if that's correct by looking at the string again.} \\
\\ 
\texttt{Row 3: ║◌│◌│C◌◌║ {-} positions: 1=◌, 2=│, 3=◌, 4=│, 5=C, 6=◌, 7=◌} \\
\\ 
\texttt{Hmm, but I said (3,3). Unless the │ characters don't count as columns? That would be strange.} \\
\\ 
\texttt{Actually, re{-}looking: maybe the grid is showing that some columns are narrower? Or perhaps I'm misunderstanding the encoding.} \\
\\ 
\texttt{Let me try yet another interpretation: the characters like │, ┤, ├, ─, etc. are part of the box{-}drawing and the actual cell contents are only the letters and ◌.} \\
\\ 
\texttt{In row 3: between the vertical borders, we have cells separated by vertical lines. So ◌│◌│C◌◌ could mean: cell1=◌, wall, cell2=◌, wall, cell3=C, cell4=◌, cell5=◌? That would be 5 cells.} \\
\\ 
\texttt{But we need 7 columns. } \\
\\ 
\texttt{Let me look at row 1 again: ║◌◌◌│◌│◌║ {-} if this is 7 cells, then │ is a cell content, not a separator.} \\
\\ 
\texttt{I think I need to accept that: columns 1{-}7 contain the characters as shown, and │, ┤, etc. are obstacles/walls that block movement. The ◌ and letters are walkable.} \\
\\ 
\texttt{So for row 3: (1,3)=◌, (2,3)=│, (3,3)=◌, (4,3)=│, (5,3)=C, (6,3)=◌, (7,3)=◌} \\
\\ 
\texttt{But then C is at (5,3), not (3,3). } \\
\\ 
\texttt{Unless... I misread my own statement? Let me check {-} I said "C at (3,3)". Maybe I made a mistake and C is actually at (5,3)?} \\
\\ 
\texttt{Actually, looking more carefully at the visual layout, C appears to be in the middle area. And (3,3) would be visually more central than (5,3).} \\
\\ 
\texttt{Wait {-} I wonder if the coordinate system is different. Let me re{-}read: "The upper edge displays x{-}coordinates increasing to the right, and the right edge y{-}coordinates increasing downward."} \\
\\ 
\texttt{So x=1 is left, y=1 is top. } \\
\\ 
\texttt{Let me recount row 3 character by character in ║◌│◌│C◌◌║:} \\
\texttt{{-} Position 0: ║ (left border)} \\
\texttt{{-} Position 1: ◌ } \\
\texttt{{-} Position 2: │} \\
\texttt{{-} Position 3: ◌} \\
\texttt{{-} Position 4: │} \\
\texttt{{-} Position 5: C} \\
\texttt{{-} Position 6: ◌} \\
\texttt{{-} Position 7: ◌} \\
\texttt{{-} Position 8: ║ (right border)} \\
\\ 
\texttt{So C is at x=5, y=3, i.e., (5,3).} \\
\\ 
\texttt{But I said (3,3). I made a mistake in my first message! } \\
\\ 
\texttt{Actually, re{-}reading {-} I said C at (3,3), L at (7,5), P at (1,7). Let me verify L and P.} \\
\\ 
\texttt{Row 5: ╟─┘◌│◌│L║ {-} wait, that's the border. The actual row 5 is ║◌├─┤◌┌─╢? No wait, let me re{-}read the grid:} \\
\\ 
\\ 
\texttt{ 1234567} \\
\texttt{╔═══╤═╤═╗} \\
\texttt{║◌◌◌│◌│◌║ 1} \\
\texttt{╟─┬─┼─┴─╢ 2} \\
\texttt{║◌│◌│C◌◌║ 3} \\
\texttt{║◌├─┤◌┌─╢ 4} \\
\texttt{╟─┘◌│◌│L║ 5} \\
\texttt{║◌◌◌└┬┴─╢ 6} \\
\texttt{║P◌◌◌│◌◌║ 7} \\
\texttt{╚════╧══╝} \\
\\ 
\\ 
\texttt{Row 5 content: ╟─┘◌│◌│L║ {-} that's the border line. The actual cells are between row 4 and row 6 borders.} \\
\\ 
\texttt{Hmm, actually looking more carefully, rows 2, 4, 5, 6 seem to have border lines. Let me re{-}interpret:} \\
\\ 
\texttt{Row 1: ║◌◌◌│◌│◌║ {-} cells} \\
\texttt{Row 2: ╟─┬─┼─┴─╢ {-} border between 1 and 2? Or is this row 2's cells?} \\
\\ 
\texttt{Actually the pattern shows: ║...║ for cell rows, and ╟...╢ for intermediate rows.} \\
\\ 
\texttt{So rows 1,3,5,7 are cell rows with ║...║} \\
\texttt{Rows 2,4,6 are intermediate rows with ╟...╢ or similar} \\
\\ 
\texttt{Wait row 4: ║◌├─┤◌┌─╢ {-} starts with ║, so it's a cell row.} \\
\texttt{Row 5: ╟─┘◌│◌│L║ {-} starts with ╟, so it's a border... but ends with ║? That seems odd.} \\
\\ 
\texttt{Actually ╟─┘◌│◌│L║ {-} this has ╟ at start and ║ at end. That's inconsistent.} \\
\\ 
\texttt{Let me re{-}read: ╟─┘◌│◌│L║ {-} maybe it's ║◌│◌│L║ with some left part? No, the string is ╟─┘◌│◌│L║.} \\
\\ 
\texttt{Hmm, looking at row 6: ║◌◌◌└┬┴─╢ {-} starts with ║, ends with ╢.} \\
\\ 
\texttt{I think there might be a pattern where some rows connect differently. Let me focus on rows with clear ║...║ pattern: 1, 3, 4, 6, 7.} \\
\\ 
\texttt{Row 1: ║◌◌◌│◌│◌║ {-} 7 cells: ◌◌◌│◌│◌} \\
\texttt{Row 3: ║◌│◌│C◌◌║ {-} 7 cells: ◌│◌│C◌◌} \\
\texttt{Row 4: ║◌├─┤◌┌─╢ {-} wait this ends with ╢ not ║} \\
\\ 
\texttt{Hmm ║◌├─┤◌┌─╢ {-} starts with ║, ends with ╢. That's 7 cells: ◌├─┤◌┌─} \\
\\ 
\texttt{Row 6: ║◌◌◌└┬┴─╢ {-} ends with ╢. 7 cells: ◌◌◌└┬┴─} \\
\\ 
\texttt{Row 7: ║P◌◌◌│◌◌║ {-} ends with ║. 7 cells: P◌◌◌│◌◌} \\
\\ 
\texttt{And L? Row 5: ╟─┘◌│◌│L║ {-} this seems to be a connector row, but has L at the end. Maybe L is in a special position?} \\
\\ 
\texttt{Actually re{-}looking: ╟─┘◌│◌│L║ {-} if we ignore the border parts, L is at position... the string is: ╟, ─, ┘, ◌, │, ◌, │, L, ║} \\
\\ 
\texttt{That's 9 characters. If the cell is at position 8 (0{-}indexed), that's x=8? But we only have 7 columns.} \\
\\ 
\texttt{I think the issue is that ╟─┘ and similar are taking up multiple character widths. In a monospace terminal, these box{-}drawing characters each take one cell.} \\
\\ 
\texttt{Let me try: ╟─┘◌│◌│L║ = ╟(1), ─(2), ┘(3), ◌(4), │(5), ◌(6), │(7), L(8), ║(9)} \\
\\ 
\texttt{So L would be at column 8? That exceeds 7.} \\
\\ 
\texttt{Unless... the row numbering on the right indicates that row 5's cells are described by this line, and L is at a valid position.} \\
\\ 
\texttt{Actually, looking at the pattern ◌│◌│L {-} if this ends the row, and we need 7 columns, then: positions 1,2,3,4,5,6,7 = something ending with ◌│◌│L?} \\
\\ 
\texttt{That would need 7 characters for positions 1{-}7. ◌│◌│L is 5 characters. With ╟─┘ prefix, that's 8 characters before ║.} \\
\\ 
\texttt{I'm going to make a different assumption: the grid uses box{-}drawing characters where some symbols like ┘ connect cells, and the actual cell contents are at specific positions.} \\
\\ 
\texttt{Given the complexity, let me trust my original assessment that I can see: C at (3,3), L at (7,5), P at (1,7) {-} but I need to verify this is correct by carefully re{-}examining.} \\
\\ 
\texttt{Actually, I think I need to step back. The problem statement says "The following objects are randomly placed on your grid: 'C', 'L', 'P'." and shows the grid. I need to parse this grid correctly.} \\
\\ 
\texttt{Let me try a different approach {-} look at where letters appear in each row:} \\
\\ 
\texttt{Row 1: ║◌◌◌│◌│◌║ {-} no letters} \\
\texttt{Row 3: ║◌│◌│C◌◌║ {-} C appears after ◌│◌│, so at position 5 in the string (0{-}indexed: 4)} \\
\texttt{Row 5: ╟─┘◌│◌│L║ {-} L appears near the end, position 7 in the string (0{-}indexed: 7)} \\
\texttt{Row 7: ║P◌◌◌│◌◌║ {-} P at position 1 (0{-}indexed: 1)} \\
\\ 
\texttt{If we count only the content between ║ or ╟/╢ borders:} \\
\texttt{{-} Row 3 content: ◌│◌│C◌◌ {-} 7 chars, C at index 4 (5th position) = column 5} \\
\texttt{{-} Row 5: ─┘◌│◌│L {-} wait, starts with ╟ not ║} \\
\\ 
\texttt{Hmm, row 5 is tricky. Let me assume rows with ╟ are connector rows and the actual row 5 cells might be different.} \\
\\ 
\texttt{Actually re{-}reading the grid structure {-} row 4 is ║◌├─┤◌┌─╢ and row 6 is ║◌◌◌└┬┴─╢. Row 5 ╟─┘◌│◌│L║ connects them.} \\
\\ 
\texttt{In box{-}drawing, ╟ connects from left, ┘ connects down and right, etc.} \\
\\ 
\texttt{I think L might be at (7,5) as I originally said, based on it appearing at the right side of row 5's display.} \\
\\ 
\texttt{Given the ambiguity, let me work with what I know: I need to find empty cells. Let me check (2,6) which the other player suggested.} \\
\\ 
\texttt{Row 6: ║◌◌◌└┬┴─╢ {-} content: ◌◌◌└┬┴─ {-} 7 chars. (1,6)=◌, (2,6)=◌, (3,6)=◌, (4,6)=└, (5,6)=┬, (6,6)=┴, (7,6)=─} \\
\\ 
\texttt{So (2,6) is ◌, empty! Good.} \\
\\ 
\texttt{(2,2): need to find row 2. The border ╟─┬─┼─┴─╢ {-} this is between rows 1 and 3, so it's not a cell row.} \\
\\ 
\texttt{Hmm, but the other player said (2,2) is empty on their board. Let me check if there's a row 2.} \\
\\ 
\texttt{Looking at the grid, y{-}coordinates go 1,2,3,4,5,6,7. So there should be row 2.} \\
\\ 
\texttt{Row 2 isn't explicitly shown with ║...║ or similar. The ╟─┬─┼─┴─╢ line is labeled as row 2 on the right side!} \\
\\ 
\texttt{So row 2 is ╟─┬─┼─┴─╢. This is a border/connector row, not a cell row with objects.} \\
\\ 
\texttt{Wait, but then where are cells for row 2? Maybe the grid is laid out differently.} \\
\\ 
\texttt{Let me re{-}interpret: the numbers 1{-}7 on the right are row labels. The lines with those labels are:} \\
\\ 
\texttt{{-} Row 1: ║◌◌◌│◌│◌║} \\
\texttt{{-} Row 2: ╟─┬─┼─┴─╢ } \\
\texttt{{-} Row 3: ║◌│◌│C◌◌║} \\
\texttt{{-} Row 4: ║◌├─┤◌┌─╢} \\
\texttt{{-} Row 5: ╟─┘◌│◌│L║} \\
\texttt{{-} Row 6: ║◌◌◌└┬┴─╢} \\
\texttt{{-} Row 7: ║P◌◌◌│◌◌║} \\
\\ 
\texttt{So rows 2 and 5 are "special" rows with ╟ instead of ║. These might still contain cells, or might be connector{-}only.} \\
\\ 
\texttt{Looking at row 5 ╟─┘◌│◌│L║ {-} it has L, so it contains at least one cell/object.} \\
\\ 
\texttt{Given ◌│◌│L at the end, and needing 7 columns, maybe positions 3{-}7 are: ◌, │, ◌, │, L? That's 5 positions. With ╟─┘ prefix (2 positions), that's 7: ┘ might be at position 2, connecting to position 3=◌.} \\
\\ 
\texttt{Actually ╟─┘ = ╟(connect left), ─(horizontal), ┘(connect down{-}right). The ┘ might indicate a corner where row 4 connects down to row 5.} \\
\\ 
\texttt{I think for practical purposes, I need to identify which (x,y) positions are empty ◌ cells. Let me focus on rows 1,3,4,6,7 which have clear ◌ patterns, and row 5 which has L.} \\
\\ 
\texttt{Row 1 ◌◌◌│◌│◌: (1,1)=◌, (2,1)=◌, (3,1)=◌, (4,1)=│, (5,1)=◌, (6,1)=│, (7,1)=◌} \\
\texttt{Row 3 ◌│◌│C◌◌: (1,3)=◌, (2,3)=│, (3,3)=◌, (4,3)=│, (5,3)=C, (6,3)=◌, (7,3)=◌} \\
\texttt{Row 4 ◌├─┤◌┌─: (1,4)=◌, (2,4)=├, (3,4)=─, (4,4)=┤, (5,4)=◌, (6,4)=┌, (7,4)=─} \\
\texttt{Row 6 ◌◌◌└┬┴─: (1,6)=◌, (2,6)=◌, (3,6)=◌, (4,6)=└, (5,6)=┬, (6,6)=┴, (7,6)=─} \\
\texttt{Row 7 P◌◌◌│◌◌: (1,7)=P, (2,7)=◌, (3,7)=◌, (4,7)=◌, (5,7)=│, (6,7)=◌, (7,7)=◌} \\
\\ 
\texttt{And row 5: need to figure out. ─┘◌│◌│L ignoring ╟ prefix? Or ┘◌│◌│L? } \\
\\ 
\texttt{If ╟─┘◌│◌│L║ and we need 7 positions between borders: maybe (1,5)=─, (2,5)=┘, (3,5)=◌, (4,5)=│, (5,5)=◌, (6,5)=│, (7,5)=L?} \\
\\ 
\texttt{That gives L at (7,5), matching my original claim!} \\
\\ 
\texttt{So summarizing my objects: C at (5,3), L at (7,5), P at (1,7).} \\
\\ 
\texttt{I made an error in my first message saying C at (3,3)! It's at (5,3).} \\
\\ 
\texttt{Now, checking empty cells from the other player's list:} \\
\texttt{{-} (4,4): I have ┤ {-} not empty} \\
\texttt{{-} (4,5): from row 5, if (4,5)=│ {-} not empty} \\
\texttt{{-} (5,4): (5,4)=◌ {-} empty!} \\
\texttt{{-} (6,4): (6,4)=┌ {-} not empty} \\
\texttt{{-} (6,5): (6,5)=│ {-} not empty} \\
\texttt{{-} (2,2): row 2 is ╟─┬─┼─┴─╢, not a cell row with ◌/letters {-} probably not valid} \\
\texttt{{-} (2,6): (2,6)=◌ {-} empty!} \\
\texttt{{-} (6,2): row 2, not a valid cell} \\
\texttt{{-} (6,6): (6,6)=┴ {-} not empty} \\
\\ 
\texttt{So common empty cells: (5,4) and (2,6)!} \\
\\ 
\texttt{Let me suggest these to the other player.} \\
\\ 
\texttt{'} \\
            }
        }
    }
    & & \\ \\

    \theutterance \stepcounter{utterance}  
    & & \multicolumn{4}{p{0.6\linewidth}}{
        \cellcolor[rgb]{0.9,0.9,0.9}{
            \makecell[{{p{\linewidth}}}]{
                \texttt{\tiny{[P1$\langle$GM]}}
                \texttt{Penalty: Your message must not contain anything before the command!} \\
\\ 
\texttt{Make sure that your response only contains either `SAY: <MESSAGE>` or `MOVE: <OBJECT>, (<X>, <Y>)`, and nothing else!} \\
\\ 
\texttt{You have collectively accumulated 3 of 8 penalties.} \\
\texttt{Please try again!} \\
            }
        }
    }
    & \\ \\

    \theutterance \stepcounter{utterance}  
    & \multicolumn{4}{p{0.6\linewidth}}{
        \cellcolor[rgb]{0.9,0.9,0.9}{
            \makecell[{{p{\linewidth}}}]{
                \texttt{\tiny{[P1$\rangle$GM]}}
                \texttt{SAY: Thanks for the list! Checking my board: (5,4) is empty ◌, (2,6) is empty ◌. But (4,4) has ┤, (6,4) has ┌, (6,5) has │, (6,6) has ┴. I need to correct my earlier message {-} my C is at (5,3), not (3,3). I misread! So: C(5,3), L(7,5), P(1,7). Can we use (5,4), (2,6) and find a third position? What about (3,1), (5,1), or (6,3)? Are those empty on your board?} \\
            }
        }
    }
    & & \\ \\

    \theutterance \stepcounter{utterance}  
    & & & \multicolumn{2}{p{0.3\linewidth}}{
        \cellcolor[rgb]{0.9,0.9,0.9}{
            \makecell[{{p{\linewidth}}}]{
                \texttt{\tiny{[GM$|$GM]}}
                \texttt{SAY: Thanks for the list! Checking my board: (5,4) is empty ◌, (2,6) is empty ◌. But (4,4) has ┤, (6,4) has ┌, (6,5) has │, (6,6) has ┴. I need to correct my earlier message {-} my C is at (5,3), not (3,3). I misread! So: C(5,3), L(7,5), P(1,7). Can we use (5,4), (2,6) and find a third position? What about (3,1), (5,1), or (6,3)? Are those empty on your board?} \\
            }
        }
    }
    & & \\ \\

    \theutterance \stepcounter{utterance}  
    & & \multicolumn{4}{p{0.6\linewidth}}{
        \cellcolor[rgb]{0.9,0.9,0.9}{
            \makecell[{{p{\linewidth}}}]{
                \texttt{\tiny{[GM$\rangle$P2]}}
                \texttt{Your message has been relayed to the other player.} \\
\\ 
\\ 
\texttt{You are currently playing round 4 of maximum 12.} \\
\texttt{You have collectively accumulated 3 of 8 penalties.} \\
\texttt{The other player sent this message:} \\
\\ 
\\ 
\texttt{"Thanks for the list! Checking my board: (5,4) is empty ◌, (2,6) is empty ◌. But (4,4) has ┤, (6,4) has ┌, (6,5) has │, (6,6) has ┴. I need to correct my earlier message {-} my C is at (5,3), not (3,3). I misread! So: C(5,3), L(7,5), P(1,7). Can we use (5,4), (2,6) and find a third position? What about (3,1), (5,1), or (6,3)? Are those empty on your board?"} \\
\\ 
\\ 
\texttt{What is your next command?} \\
            }
        }
    }
    & \\ \\

    \theutterance \stepcounter{utterance}  
    & & & \multicolumn{4}{p{0.6\linewidth}}{
        \cellcolor[rgb]{0.9,0.9,0.9}{
            \makecell[{{p{\linewidth}}}]{
                \texttt{\tiny{[GM$\langle$P2]}}
                \texttt{SAY: On my board: (3,1) is ◌, (5,1) is │, (6,3) is ◌. So (3,1) and (6,3) work! Let's use C→(5,4), L→(3,1), P→(6,3). Does this work for you?} \\
            }
        }
    }
     \\ \\

    \theutterance \stepcounter{utterance}  
    & & & \multicolumn{2}{p{0.3\linewidth}}{
        \cellcolor[rgb]{0.9,0.9,0.9}{
            \makecell[{{p{\linewidth}}}]{
                \texttt{\tiny{[GM$|$GM]}}
                \texttt{SAY: On my board: (3,1) is ◌, (5,1) is │, (6,3) is ◌. So (3,1) and (6,3) work! Let's use C→(5,4), L→(3,1), P→(6,3). Does this work for you?} \\
            }
        }
    }
    & & \\ \\

    \theutterance \stepcounter{utterance}  
    & & \multicolumn{4}{p{0.6\linewidth}}{
        \cellcolor[rgb]{0.9,0.9,0.9}{
            \makecell[{{p{\linewidth}}}]{
                \texttt{\tiny{[P1$\langle$GM]}}
                \texttt{Your message has been relayed to the other player.} \\
\\ 
\\ 
\texttt{You are currently playing round 4 of maximum 12.} \\
\texttt{You have collectively accumulated 3 of 8 penalties.} \\
\texttt{The other player sent this message:} \\
\\ 
\\ 
\texttt{"On my board: (3,1) is ◌, (5,1) is │, (6,3) is ◌. So (3,1) and (6,3) work! Let's use C→(5,4), L→(3,1), P→(6,3). Does this work for you?"} \\
\\ 
\\ 
\texttt{What is your next command?} \\
            }
        }
    }
    & \\ \\

    \theutterance \stepcounter{utterance}  
    & \multicolumn{4}{p{0.6\linewidth}}{
        \cellcolor[rgb]{0.9,0.9,0.9}{
            \makecell[{{p{\linewidth}}}]{
                \texttt{\tiny{[P1$\rangle$GM]}}
                \texttt{SAY: Perfect! Let me verify: C→(5,4) is ◌, L→(3,1) is ◌, P→(6,3) is ◌. All empty on my board! Let's do this order: C first, then L, then P. I'll move my C from (5,3) to (5,4) now.} \\
            }
        }
    }
    & & \\ \\

    \theutterance \stepcounter{utterance}  
    & & & \multicolumn{2}{p{0.3\linewidth}}{
        \cellcolor[rgb]{0.9,0.9,0.9}{
            \makecell[{{p{\linewidth}}}]{
                \texttt{\tiny{[GM$|$GM]}}
                \texttt{SAY: Perfect! Let me verify: C→(5,4) is ◌, L→(3,1) is ◌, P→(6,3) is ◌. All empty on my board! Let's do this order: C first, then L, then P. I'll move my C from (5,3) to (5,4) now.} \\
            }
        }
    }
    & & \\ \\

    \theutterance \stepcounter{utterance}  
    & & \multicolumn{4}{p{0.6\linewidth}}{
        \cellcolor[rgb]{0.9,0.9,0.9}{
            \makecell[{{p{\linewidth}}}]{
                \texttt{\tiny{[GM$\rangle$P2]}}
                \texttt{Your message has been relayed to the other player.} \\
\\ 
\\ 
\texttt{You are currently playing round 5 of maximum 12.} \\
\texttt{You have collectively accumulated 3 of 8 penalties.} \\
\texttt{The other player sent this message:} \\
\\ 
\\ 
\texttt{"Perfect! Let me verify: C→(5,4) is ◌, L→(3,1) is ◌, P→(6,3) is ◌. All empty on my board! Let's do this order: C first, then L, then P. I'll move my C from (5,3) to (5,4) now."} \\
\\ 
\\ 
\texttt{What is your next command?} \\
            }
        }
    }
    & \\ \\

    \theutterance \stepcounter{utterance}  
    & & & \multicolumn{4}{p{0.6\linewidth}}{
        \cellcolor[rgb]{0.9,0.9,0.9}{
            \makecell[{{p{\linewidth}}}]{
                \texttt{\tiny{[GM$\langle$P2]}}
                \texttt{MOVE: C, (5,4)} \\
            }
        }
    }
     \\ \\

    \theutterance \stepcounter{utterance}  
    & & & \multicolumn{2}{p{0.3\linewidth}}{
        \cellcolor[rgb]{0.9,0.9,0.9}{
            \makecell[{{p{\linewidth}}}]{
                \texttt{\tiny{[GM$|$GM]}}
                \texttt{MOVE: C, (5,4)} \\
            }
        }
    }
    & & \\ \\

    \theutterance \stepcounter{utterance}  
    & & & \multicolumn{2}{p{0.3\linewidth}}{
        \cellcolor[rgb]{0.9,0.9,0.9}{
            \makecell[{{p{\linewidth}}}]{
                \texttt{\tiny{[GM$|$GM]}}
                \texttt{C moved to (5, 4)} \\
            }
        }
    }
    & & \\ \\

    \theutterance \stepcounter{utterance}  
    & & \multicolumn{4}{p{0.6\linewidth}}{
        \cellcolor[rgb]{0.9,0.9,0.9}{
            \makecell[{{p{\linewidth}}}]{
                \texttt{\tiny{[P1$\langle$GM]}}
                \texttt{Your message has been relayed to the other player.} \\
\\ 
\\ 
\texttt{You are currently playing round 5 of maximum 12.} \\
\texttt{You have collectively accumulated 3 of 8 penalties.} \\
\texttt{The other player moved an object on their grid.} \\
\\ 
\\ 
\texttt{What is your next command?} \\
            }
        }
    }
    & \\ \\

    \theutterance \stepcounter{utterance}  
    & \multicolumn{4}{p{0.6\linewidth}}{
        \cellcolor[rgb]{0.9,0.9,0.9}{
            \makecell[{{p{\linewidth}}}]{
                \texttt{\tiny{[P1$\rangle$GM]}}
                \texttt{MOVE: C, (5, 4)} \\
            }
        }
    }
    & & \\ \\

    \theutterance \stepcounter{utterance}  
    & & & \multicolumn{2}{p{0.3\linewidth}}{
        \cellcolor[rgb]{0.9,0.9,0.9}{
            \makecell[{{p{\linewidth}}}]{
                \texttt{\tiny{[GM$|$GM]}}
                \texttt{MOVE: C, (5, 4)} \\
            }
        }
    }
    & & \\ \\

    \theutterance \stepcounter{utterance}  
    & & & \multicolumn{2}{p{0.3\linewidth}}{
        \cellcolor[rgb]{0.9,0.9,0.9}{
            \makecell[{{p{\linewidth}}}]{
                \texttt{\tiny{[GM$|$GM]}}
                \texttt{C moved to (5, 4)} \\
            }
        }
    }
    & & \\ \\

    \theutterance \stepcounter{utterance}  
    & & \multicolumn{4}{p{0.6\linewidth}}{
        \cellcolor[rgb]{0.9,0.9,0.9}{
            \makecell[{{p{\linewidth}}}]{
                \texttt{\tiny{[GM$\rangle$P2]}}
                \texttt{Moved 'C' to (5,4) successfully. Your updated grid looks like this:} \\
\texttt{```} \\
\texttt{ 1234567} \\
\texttt{╔═══╤═╤═╗} \\
\texttt{║◌◌◌│◌│◌║ 1} \\
\texttt{╟─┬─┼─┴─╢ 2} \\
\texttt{║◌│L│◌P◌║ 3} \\
\texttt{║◌├─┤C┌─╢ 4} \\
\texttt{╟─┘◌│◌│◌║ 5} \\
\texttt{║◌◌◌└┬┴─╢ 6} \\
\texttt{║◌◌◌◌│◌◌║ 7} \\
\texttt{╚════╧══╝} \\
\texttt{```} \\
\texttt{You are currently playing round 6 of maximum 12.} \\
\texttt{You have collectively accumulated 3 of 8 penalties.} \\
\texttt{The other player moved an object on their grid.} \\
\\ 
\\ 
\texttt{What is your next command?} \\
            }
        }
    }
    & \\ \\

    \theutterance \stepcounter{utterance}  
    & & & \multicolumn{4}{p{0.6\linewidth}}{
        \cellcolor[rgb]{0.9,0.9,0.9}{
            \makecell[{{p{\linewidth}}}]{
                \texttt{\tiny{[GM$\langle$P2]}}
                \texttt{SAY: C moved successfully! Now moving L to (3,1).} \\
            }
        }
    }
     \\ \\

    \theutterance \stepcounter{utterance}  
    & & & \multicolumn{2}{p{0.3\linewidth}}{
        \cellcolor[rgb]{0.9,0.9,0.9}{
            \makecell[{{p{\linewidth}}}]{
                \texttt{\tiny{[GM$|$GM]}}
                \texttt{SAY: C moved successfully! Now moving L to (3,1).} \\
            }
        }
    }
    & & \\ \\

    \theutterance \stepcounter{utterance}  
    & & \multicolumn{4}{p{0.6\linewidth}}{
        \cellcolor[rgb]{0.9,0.9,0.9}{
            \makecell[{{p{\linewidth}}}]{
                \texttt{\tiny{[P1$\langle$GM]}}
                \texttt{Moved 'C' to (5,4) successfully. Your updated grid looks like this:} \\
\texttt{```} \\
\texttt{ 1234567} \\
\texttt{╔═══╤═╤═╗} \\
\texttt{║◌◌◌│◌│◌║ 1} \\
\texttt{╟─┬─┼─┴─╢ 2} \\
\texttt{║◌│◌│◌◌◌║ 3} \\
\texttt{║◌├─┤C┌─╢ 4} \\
\texttt{╟─┘◌│◌│L║ 5} \\
\texttt{║◌◌◌└┬┴─╢ 6} \\
\texttt{║P◌◌◌│◌◌║ 7} \\
\texttt{╚════╧══╝} \\
\texttt{```} \\
\texttt{You are currently playing round 6 of maximum 12.} \\
\texttt{You have collectively accumulated 3 of 8 penalties.} \\
\texttt{The other player sent this message:} \\
\\ 
\\ 
\texttt{"C moved successfully! Now moving L to (3,1)."} \\
\\ 
\\ 
\texttt{What is your next command?} \\
            }
        }
    }
    & \\ \\

    \theutterance \stepcounter{utterance}  
    & \multicolumn{4}{p{0.6\linewidth}}{
        \cellcolor[rgb]{0.9,0.9,0.9}{
            \makecell[{{p{\linewidth}}}]{
                \texttt{\tiny{[P1$\rangle$GM]}}
                \texttt{MOVE: L, (3, 1)} \\
            }
        }
    }
    & & \\ \\

    \theutterance \stepcounter{utterance}  
    & & & \multicolumn{2}{p{0.3\linewidth}}{
        \cellcolor[rgb]{0.9,0.9,0.9}{
            \makecell[{{p{\linewidth}}}]{
                \texttt{\tiny{[GM$|$GM]}}
                \texttt{MOVE: L, (3, 1)} \\
            }
        }
    }
    & & \\ \\

    \theutterance \stepcounter{utterance}  
    & & & \multicolumn{2}{p{0.3\linewidth}}{
        \cellcolor[rgb]{0.9,0.9,0.9}{
            \makecell[{{p{\linewidth}}}]{
                \texttt{\tiny{[GM$|$GM]}}
                \texttt{L moved to (3, 1)} \\
            }
        }
    }
    & & \\ \\

    \theutterance \stepcounter{utterance}  
    & & \multicolumn{4}{p{0.6\linewidth}}{
        \cellcolor[rgb]{0.9,0.9,0.9}{
            \makecell[{{p{\linewidth}}}]{
                \texttt{\tiny{[GM$\rangle$P2]}}
                \texttt{Your message has been relayed to the other player.} \\
\\ 
\\ 
\texttt{You are currently playing round 7 of maximum 12.} \\
\texttt{You have collectively accumulated 3 of 8 penalties.} \\
\texttt{The other player moved an object on their grid.} \\
\\ 
\\ 
\texttt{What is your next command?} \\
            }
        }
    }
    & \\ \\

    \theutterance \stepcounter{utterance}  
    & & & \multicolumn{4}{p{0.6\linewidth}}{
        \cellcolor[rgb]{0.9,0.9,0.9}{
            \makecell[{{p{\linewidth}}}]{
                \texttt{\tiny{[GM$\langle$P2]}}
                \texttt{MOVE: L, (3,1)} \\
            }
        }
    }
     \\ \\

    \theutterance \stepcounter{utterance}  
    & & & \multicolumn{2}{p{0.3\linewidth}}{
        \cellcolor[rgb]{0.9,0.9,0.9}{
            \makecell[{{p{\linewidth}}}]{
                \texttt{\tiny{[GM$|$GM]}}
                \texttt{MOVE: L, (3,1)} \\
            }
        }
    }
    & & \\ \\

    \theutterance \stepcounter{utterance}  
    & & & \multicolumn{2}{p{0.3\linewidth}}{
        \cellcolor[rgb]{0.9,0.9,0.9}{
            \makecell[{{p{\linewidth}}}]{
                \texttt{\tiny{[GM$|$GM]}}
                \texttt{L moved to (3, 1)} \\
            }
        }
    }
    & & \\ \\

    \theutterance \stepcounter{utterance}  
    & & \multicolumn{4}{p{0.6\linewidth}}{
        \cellcolor[rgb]{0.9,0.9,0.9}{
            \makecell[{{p{\linewidth}}}]{
                \texttt{\tiny{[P1$\langle$GM]}}
                \texttt{Moved 'L' to (3,1) successfully. Your updated grid looks like this:} \\
\texttt{```} \\
\texttt{ 1234567} \\
\texttt{╔═══╤═╤═╗} \\
\texttt{║◌◌L│◌│◌║ 1} \\
\texttt{╟─┬─┼─┴─╢ 2} \\
\texttt{║◌│◌│◌◌◌║ 3} \\
\texttt{║◌├─┤C┌─╢ 4} \\
\texttt{╟─┘◌│◌│◌║ 5} \\
\texttt{║◌◌◌└┬┴─╢ 6} \\
\texttt{║P◌◌◌│◌◌║ 7} \\
\texttt{╚════╧══╝} \\
\texttt{```} \\
\texttt{You are currently playing round 7 of maximum 12.} \\
\texttt{You have collectively accumulated 3 of 8 penalties.} \\
\texttt{The other player moved an object on their grid.} \\
\\ 
\\ 
\texttt{What is your next command?} \\
            }
        }
    }
    & \\ \\

    \theutterance \stepcounter{utterance}  
    & \multicolumn{4}{p{0.6\linewidth}}{
        \cellcolor[rgb]{0.9,0.9,0.9}{
            \makecell[{{p{\linewidth}}}]{
                \texttt{\tiny{[P1$\rangle$GM]}}
                \texttt{MOVE: P, (6, 3)} \\
            }
        }
    }
    & & \\ \\

    \theutterance \stepcounter{utterance}  
    & & & \multicolumn{2}{p{0.3\linewidth}}{
        \cellcolor[rgb]{0.9,0.9,0.9}{
            \makecell[{{p{\linewidth}}}]{
                \texttt{\tiny{[GM$|$GM]}}
                \texttt{MOVE: P, (6, 3)} \\
            }
        }
    }
    & & \\ \\

    \theutterance \stepcounter{utterance}  
    & & & \multicolumn{2}{p{0.3\linewidth}}{
        \cellcolor[rgb]{0.9,0.9,0.9}{
            \makecell[{{p{\linewidth}}}]{
                \texttt{\tiny{[GM$|$GM]}}
                \texttt{P moved to (6, 3)} \\
            }
        }
    }
    & & \\ \\

    \theutterance \stepcounter{utterance}  
    & & \multicolumn{4}{p{0.6\linewidth}}{
        \cellcolor[rgb]{0.9,0.9,0.9}{
            \makecell[{{p{\linewidth}}}]{
                \texttt{\tiny{[GM$\rangle$P2]}}
                \texttt{Moved 'L' to (3,1) successfully. Your updated grid looks like this:} \\
\texttt{```} \\
\texttt{ 1234567} \\
\texttt{╔═══╤═╤═╗} \\
\texttt{║◌◌L│◌│◌║ 1} \\
\texttt{╟─┬─┼─┴─╢ 2} \\
\texttt{║◌│◌│◌P◌║ 3} \\
\texttt{║◌├─┤C┌─╢ 4} \\
\texttt{╟─┘◌│◌│◌║ 5} \\
\texttt{║◌◌◌└┬┴─╢ 6} \\
\texttt{║◌◌◌◌│◌◌║ 7} \\
\texttt{╚════╧══╝} \\
\texttt{```} \\
\texttt{You are currently playing round 8 of maximum 12.} \\
\texttt{You have collectively accumulated 3 of 8 penalties.} \\
\texttt{The other player moved an object on their grid.} \\
\\ 
\\ 
\texttt{What is your next command?} \\
            }
        }
    }
    & \\ \\

    \theutterance \stepcounter{utterance}  
    & & & \multicolumn{4}{p{0.6\linewidth}}{
        \cellcolor[rgb]{0.9,0.9,0.9}{
            \makecell[{{p{\linewidth}}}]{
                \texttt{\tiny{[GM$\langle$P2]}}
                \texttt{SAY: L moved successfully! Now moving P to (6,3).} \\
            }
        }
    }
     \\ \\

    \theutterance \stepcounter{utterance}  
    & & & \multicolumn{2}{p{0.3\linewidth}}{
        \cellcolor[rgb]{0.9,0.9,0.9}{
            \makecell[{{p{\linewidth}}}]{
                \texttt{\tiny{[GM$|$GM]}}
                \texttt{SAY: L moved successfully! Now moving P to (6,3).} \\
            }
        }
    }
    & & \\ \\

    \theutterance \stepcounter{utterance}  
    & & \multicolumn{4}{p{0.6\linewidth}}{
        \cellcolor[rgb]{0.9,0.9,0.9}{
            \makecell[{{p{\linewidth}}}]{
                \texttt{\tiny{[P1$\langle$GM]}}
                \texttt{Moved 'P' to (6,3) successfully. Your updated grid looks like this:} \\
\texttt{```} \\
\texttt{ 1234567} \\
\texttt{╔═══╤═╤═╗} \\
\texttt{║◌◌L│◌│◌║ 1} \\
\texttt{╟─┬─┼─┴─╢ 2} \\
\texttt{║◌│◌│◌P◌║ 3} \\
\texttt{║◌├─┤C┌─╢ 4} \\
\texttt{╟─┘◌│◌│◌║ 5} \\
\texttt{║◌◌◌└┬┴─╢ 6} \\
\texttt{║◌◌◌◌│◌◌║ 7} \\
\texttt{╚════╧══╝} \\
\texttt{```} \\
\texttt{You are currently playing round 8 of maximum 12.} \\
\texttt{You have collectively accumulated 3 of 8 penalties.} \\
\texttt{The other player sent this message:} \\
\\ 
\\ 
\texttt{"L moved successfully! Now moving P to (6,3)."} \\
\\ 
\\ 
\texttt{What is your next command?} \\
            }
        }
    }
    & \\ \\

    \theutterance \stepcounter{utterance}  
    & \multicolumn{4}{p{0.6\linewidth}}{
        \cellcolor[rgb]{0.9,0.9,0.9}{
            \makecell[{{p{\linewidth}}}]{
                \texttt{\tiny{[P1$\rangle$GM]}}
                \texttt{SAY: All my objects are now in position: C at (5,4), L at (3,1), P at (6,3). Are your objects also at these coordinates?} \\
            }
        }
    }
    & & \\ \\

    \theutterance \stepcounter{utterance}  
    & & & \multicolumn{2}{p{0.3\linewidth}}{
        \cellcolor[rgb]{0.9,0.9,0.9}{
            \makecell[{{p{\linewidth}}}]{
                \texttt{\tiny{[GM$|$GM]}}
                \texttt{SAY: All my objects are now in position: C at (5,4), L at (3,1), P at (6,3). Are your objects also at these coordinates?} \\
            }
        }
    }
    & & \\ \\

    \theutterance \stepcounter{utterance}  
    & & \multicolumn{4}{p{0.6\linewidth}}{
        \cellcolor[rgb]{0.9,0.9,0.9}{
            \makecell[{{p{\linewidth}}}]{
                \texttt{\tiny{[GM$\rangle$P2]}}
                \texttt{Your message has been relayed to the other player.} \\
\\ 
\\ 
\texttt{You are currently playing round 9 of maximum 12.} \\
\texttt{You have collectively accumulated 3 of 8 penalties.} \\
\texttt{The other player sent this message:} \\
\\ 
\\ 
\texttt{"All my objects are now in position: C at (5,4), L at (3,1), P at (6,3). Are your objects also at these coordinates?"} \\
\\ 
\\ 
\texttt{What is your next command?} \\
            }
        }
    }
    & \\ \\

    \theutterance \stepcounter{utterance}  
    & & & \multicolumn{4}{p{0.6\linewidth}}{
        \cellcolor[rgb]{0.9,0.9,0.9}{
            \makecell[{{p{\linewidth}}}]{
                \texttt{\tiny{[GM$\langle$P2]}}
                \texttt{MOVE: P, (6,3)} \\
            }
        }
    }
     \\ \\

    \theutterance \stepcounter{utterance}  
    & & & \multicolumn{2}{p{0.3\linewidth}}{
        \cellcolor[rgb]{0.9,0.9,0.9}{
            \makecell[{{p{\linewidth}}}]{
                \texttt{\tiny{[GM$|$GM]}}
                \texttt{MOVE: P, (6,3)} \\
            }
        }
    }
    & & \\ \\

    \theutterance \stepcounter{utterance}  
    & & & \multicolumn{2}{p{0.3\linewidth}}{
        \cellcolor[rgb]{0.9,0.9,0.9}{
            \makecell[{{p{\linewidth}}}]{
                \texttt{\tiny{[GM$|$GM]}}
                \texttt{Penalty: (6,3) is not empty, but contains 'P'.} \\
\\ 
\texttt{You have collectively accumulated 4 of 8 penalties.} \\
\texttt{Please try again!} \\
            }
        }
    }
    & & \\ \\

    \theutterance \stepcounter{utterance}  
    & & \multicolumn{4}{p{0.6\linewidth}}{
        \cellcolor[rgb]{0.9,0.9,0.9}{
            \makecell[{{p{\linewidth}}}]{
                \texttt{\tiny{[GM$\rangle$P2]}}
                \texttt{Penalty: (6,3) is not empty, but contains 'P'.} \\
\\ 
\texttt{You have collectively accumulated 4 of 8 penalties.} \\
\texttt{Please try again!} \\
            }
        }
    }
    & \\ \\

    \theutterance \stepcounter{utterance}  
    & & & \multicolumn{4}{p{0.6\linewidth}}{
        \cellcolor[rgb]{0.9,0.9,0.9}{
            \makecell[{{p{\linewidth}}}]{
                \texttt{\tiny{[GM$\langle$P2]}}
                \texttt{SAY: finished?} \\
            }
        }
    }
     \\ \\

    \theutterance \stepcounter{utterance}  
    & & & \multicolumn{2}{p{0.3\linewidth}}{
        \cellcolor[rgb]{0.9,0.9,0.9}{
            \makecell[{{p{\linewidth}}}]{
                \texttt{\tiny{[GM$|$GM]}}
                \texttt{SAY: finished?} \\
            }
        }
    }
    & & \\ \\

    \theutterance \stepcounter{utterance}  
    & & \multicolumn{4}{p{0.6\linewidth}}{
        \cellcolor[rgb]{0.9,0.9,0.9}{
            \makecell[{{p{\linewidth}}}]{
                \texttt{\tiny{[P1$\langle$GM]}}
                \texttt{Your message has been relayed to the other player.} \\
\\ 
\\ 
\texttt{You are currently playing round 9 of maximum 12.} \\
\texttt{You have collectively accumulated 4 of 8 penalties.} \\
\texttt{The other player sent this message:} \\
\\ 
\\ 
\texttt{"finished?"} \\
\\ 
\\ 
\texttt{What is your next command?} \\
            }
        }
    }
    & \\ \\

    \theutterance \stepcounter{utterance}  
    & \multicolumn{4}{p{0.6\linewidth}}{
        \cellcolor[rgb]{0.9,0.9,0.9}{
            \makecell[{{p{\linewidth}}}]{
                \texttt{\tiny{[P1$\rangle$GM]}}
                \texttt{SAY: finished!} \\
            }
        }
    }
    & & \\ \\

    \theutterance \stepcounter{utterance}  
    & & & \multicolumn{2}{p{0.3\linewidth}}{
        \cellcolor[rgb]{0.9,0.9,0.9}{
            \makecell[{{p{\linewidth}}}]{
                \texttt{\tiny{[GM$|$GM]}}
                \texttt{SAY: finished!} \\
            }
        }
    }
    & & \\ \\

    \theutterance \stepcounter{utterance}  
    & & & \multicolumn{2}{p{0.3\linewidth}}{
        \cellcolor[rgb]{0.9,0.9,0.9}{
            \makecell[{{p{\linewidth}}}]{
                \texttt{\tiny{[GM$|$GM]}}
                \texttt{true} \\
            }
        }
    }
    & & \\ \\

    \theutterance \stepcounter{utterance}  
    & & & \multicolumn{2}{p{0.3\linewidth}}{
        \cellcolor[rgb]{0.9,0.9,0.9}{
            \makecell[{{p{\linewidth}}}]{
                \texttt{\tiny{[GM$|$GM]}}
                \texttt{Initial states:} \\
\texttt{```} \\
\texttt{Player 1:} \\
\texttt{ 1234567} \\
\texttt{╔═══╤═╤═╗} \\
\texttt{║◌◌◌│◌│◌║ 1} \\
\texttt{╟─┬─┼─┴─╢ 2} \\
\texttt{║◌│◌│C◌◌║ 3} \\
\texttt{║◌├─┤◌┌─╢ 4} \\
\texttt{╟─┘◌│◌│L║ 5} \\
\texttt{║◌◌◌└┬┴─╢ 6} \\
\texttt{║P◌◌◌│◌◌║ 7} \\
\texttt{╚════╧══╝} \\
\\ 
\\ 
\texttt{Player 2:} \\
\texttt{ 1234567} \\
\texttt{╔═══╤═╤═╗} \\
\texttt{║◌◌◌│◌│◌║ 1} \\
\texttt{╟─┬─┼─┴─╢ 2} \\
\texttt{║◌│L│◌P◌║ 3} \\
\texttt{║◌├─┤◌┌─╢ 4} \\
\texttt{╟─┘◌│◌│◌║ 5} \\
\texttt{║◌◌◌└┬┴─╢ 6} \\
\texttt{║◌C◌◌│◌◌║ 7} \\
\texttt{╚════╧══╝} \\
\\ 
\texttt{```} \\
            }
        }
    }
    & & \\ \\

    \theutterance \stepcounter{utterance}  
    & & & \multicolumn{2}{p{0.3\linewidth}}{
        \cellcolor[rgb]{0.9,0.9,0.9}{
            \makecell[{{p{\linewidth}}}]{
                \texttt{\tiny{[GM$|$GM]}}
                \texttt{End states:} \\
\texttt{```} \\
\texttt{Player 1:} \\
\texttt{ 1234567} \\
\texttt{╔═══╤═╤═╗} \\
\texttt{║◌◌L│◌│◌║ 1} \\
\texttt{╟─┬─┼─┴─╢ 2} \\
\texttt{║◌│◌│◌P◌║ 3} \\
\texttt{║◌├─┤C┌─╢ 4} \\
\texttt{╟─┘◌│◌│◌║ 5} \\
\texttt{║◌◌◌└┬┴─╢ 6} \\
\texttt{║◌◌◌◌│◌◌║ 7} \\
\texttt{╚════╧══╝} \\
\\ 
\\ 
\texttt{Player 2:} \\
\texttt{ 1234567} \\
\texttt{╔═══╤═╤═╗} \\
\texttt{║◌◌L│◌│◌║ 1} \\
\texttt{╟─┬─┼─┴─╢ 2} \\
\texttt{║◌│◌│◌P◌║ 3} \\
\texttt{║◌├─┤C┌─╢ 4} \\
\texttt{╟─┘◌│◌│◌║ 5} \\
\texttt{║◌◌◌└┬┴─╢ 6} \\
\texttt{║◌◌◌◌│◌◌║ 7} \\
\texttt{╚════╧══╝} \\
\\ 
\texttt{```} \\
            }
        }
    }
    & & \\ \\

    \theutterance \stepcounter{utterance}  
    & & & \multicolumn{2}{p{0.3\linewidth}}{
        \cellcolor[rgb]{0.9,0.9,0.9}{
            \makecell[{{p{\linewidth}}}]{
                \texttt{\tiny{[GM$|$GM]}}
                \texttt{* success: True} \\
\texttt{* lose: False} \\
\texttt{* aborted: False} \\
\texttt{{-}{-}{-}{-}{-}{-}{-}} \\
\texttt{* Shifts: 2.00} \\
\texttt{* Max Shifts: 4.00} \\
\texttt{* Min Shifts: 2.00} \\
\texttt{* End Distance Sum: 0.00} \\
\texttt{* Init Distance Sum: 15.88} \\
\texttt{* Expected Distance Sum: 12.57} \\
\texttt{* Penalties: 4.00} \\
\texttt{* Max Penalties: 8.00} \\
\texttt{* Rounds: 9.00} \\
\texttt{* Max Rounds: 12.00} \\
\texttt{* Object Count: 3.00} \\
            }
        }
    }
    & & \\ \\

    \theutterance \stepcounter{utterance}  
    & & & \multicolumn{2}{p{0.3\linewidth}}{
        \cellcolor[rgb]{0.9,0.9,0.9}{
            \makecell[{{p{\linewidth}}}]{
                \texttt{\tiny{[GM$|$GM]}}
                \texttt{Initial Grids:} \\
\texttt{```} \\
\texttt{Player 1:     Player 2:} \\
\texttt{ 1234567       1234567} \\
\texttt{╔═══╤═╤═╗     ╔═══╤═╤═╗} \\
\texttt{║◌◌◌│◌│◌║ 1   ║◌◌◌│◌│◌║ 1} \\
\texttt{╟─┬─┼─┴─╢ 2   ╟─┬─┼─┴─╢ 2} \\
\texttt{║◌│◌│C◌◌║ 3   ║◌│L│◌P◌║ 3} \\
\texttt{║◌├─┤◌┌─╢ 4   ║◌├─┤◌┌─╢ 4} \\
\texttt{╟─┘◌│◌│L║ 5   ╟─┘◌│◌│◌║ 5} \\
\texttt{║◌◌◌└┬┴─╢ 6   ║◌◌◌└┬┴─╢ 6} \\
\texttt{║P◌◌◌│◌◌║ 7   ║◌C◌◌│◌◌║ 7} \\
\texttt{╚════╧══╝     ╚════╧══╝} \\
\texttt{```} \\
\\ 
\texttt{Final Grids:} \\
\texttt{```} \\
\texttt{Player 1:     Player 2:} \\
\texttt{ 1234567       1234567} \\
\texttt{╔═══╤═╤═╗     ╔═══╤═╤═╗} \\
\texttt{║◌◌L│◌│◌║ 1   ║◌◌L│◌│◌║ 1} \\
\texttt{╟─┬─┼─┴─╢ 2   ╟─┬─┼─┴─╢ 2} \\
\texttt{║◌│◌│◌P◌║ 3   ║◌│◌│◌P◌║ 3} \\
\texttt{║◌├─┤C┌─╢ 4   ║◌├─┤C┌─╢ 4} \\
\texttt{╟─┘◌│◌│◌║ 5   ╟─┘◌│◌│◌║ 5} \\
\texttt{║◌◌◌└┬┴─╢ 6   ║◌◌◌└┬┴─╢ 6} \\
\texttt{║◌◌◌◌│◌◌║ 7   ║◌◌◌◌│◌◌║ 7} \\
\texttt{╚════╧══╝     ╚════╧══╝} \\
\texttt{```} \\
            }
        }
    }
    & & \\ \\

\end{supertabular}
}

\end{document}
